\documentclass[justified,openright]{tufte-book}
\usepackage{amsmath,amssymb,amsfonts}

% ------------------------
% Typography
% ------------------------
\usepackage[protrusion=false]{microtype}
\usepackage{fontspec}
\setmainfont{TeX Gyre Pagella}
\setsansfont[Scale=0.92]{TeX Gyre Heros}
\setmonofont[Scale=0.82]{TeX Gyre Cursor}
% ------------------------
% Graphics and tables
% ------------------------
\usepackage{graphicx}
\usepackage{booktabs}
\usepackage{longtable}
\usepackage{array}
% ------------------------
% Colors
% ------------------------
\usepackage[table,xcdraw,svgnames]{xcolor}

\definecolor{DarkGray}{HTML}{444444}
\definecolor{CalloutBlue}{HTML}{DCEEFF}
\definecolor{CalloutBorder}{HTML}{5B8BD9}
\definecolor{BrandBlue}{RGB}{18,65,124}
\definecolor{BrandAccent}{RGB}{0,153,204}
\definecolor{LightGray}{RGB}{245,246,248}
\definecolor{MidGray}{RGB}{175,180,190}
\definecolor{DarkGray}{RGB}{55,60,70}
\definecolor{WarningAmber}{RGB}{230,100,0}
\definecolor{WarningBg}{RGB}{255,248,235}
\definecolor{CautionYellow}{RGB}{180,140,0}
\definecolor{CautionBg}{RGB}{255,252,230}
\definecolor{TipGreen}{RGB}{36,110,44}
\definecolor{TipBg}{RGB}{237,248,239}
\definecolor{NoteBg}{RGB}{235,243,255}
\definecolor{NoteBlue}{RGB}{18,90,160}
\definecolor{DangerRed}{RGB}{185,28,28}
\definecolor{DangerBg}{RGB}{255,237,237}
\definecolor{TableHeader}{RGB}{18,65,124}
\definecolor{TableRowAlt}{RGB}{235,241,252}
\definecolor{MarginRule}{RGB}{210,215,225}
\definecolor{SpineAccent}{RGB}{0,153,204}
% ------------------------
% Utility packages
% ------------------------
\usepackage{etoolbox}
\usepackage{xparse}
\usepackage{ifthen}

% ------------------------
% Margin notes (native to tufte-book)
% ------------------------
\setlength{\marginparwidth}{2.2in}

\newcommand{\mn}[1]{%
  \marginnote{\raggedright\footnotesize\color{DarkGray}#1}
}

% ------------------------
% Margin callout box
% ------------------------
\NewDocumentCommand{\margincallout}{m}{%
  \marginnote{%
    \setlength{\fboxsep}{6pt}%
    \colorbox{CalloutBlue}{%
      \begin{minipage}{\marginparwidth}
      \footnotesize #1
      \end{minipage}
    }
  }
}

% ------------------------
% Paragraph formatting
% ------------------------
\setlength{\parindent}{1em}
\setlength{\parskip}{0pt}

% ------------------------
% Section spacing tweaks
% ------------------------
\usepackage{titlesec}

\titlespacing*{\section}
  {0pt}{2.0ex plus .2ex}{1.0ex}

\titlespacing*{\subsection}
  {0pt}{1.5ex plus .2ex}{0.8ex}

% ------------------------
% Hyperlinks
% ------------------------
\usepackage{hyperref}
\usepackage{tabularray}
\UseTblrLibrary{booktabs}

\usepackage{soul}
\AfterEndEnvironment{longtblr}{\soulsuspend} 
%\makeatletter
%\let\MakeTextLowercase\MakeLowercase
%\makeatother

\makeatletter
% Provide safe no-op fallbacks (tabularray may invoke these)
\providecommand{\soulsuspend}{}
\providecommand{\soulsuspend*}{}
\providecommand{\soulsuspend@}{}
\makeatother

% Robust checkbox macro (do NOT make it expand to itself)
\DeclareRobustCommand{\checkbox}{\(\square\)}

% Register commonly used macros so the `soul` reconstruction succeeds.
% Keep registration after the macro is defined.
\soulregister\checkbox{0}
\soulregister\textbf{1}
\soulregister\emph{1}
\soulregister\href{2}
\soulregister\label{0}
\soulregister\ref{0}
\soulregister\pageref{0}
\soulregister\footnote{1}
\soulregister\cite{0}

% -- Suppress cosmetic warnings
\usepackage{silence}
\WarningFilter{latex}{Marginpar on page}
\WarningFilter{bibentry}{You have used}
\WarningFilter{rerunfilecheck}{File}


\usepackage{textcase} % tufte loads it, but this keeps it explicit

\makeatletter
% 1) Put *pure* titles into marks (no case-change at mark creation)
\renewcommand{\chaptermark}[1]{\markboth{#1}{}}
\renewcommand{\sectionmark}[1]{\markright{#1}}

% 2) Prevent 'soul' letterspacing macros from entering marks
\pretocmd\markboth{\let\allcapsspacing\relax \let\smallcapsspacing\relax}{}{}
\pretocmd\markright{\let\allcapsspacing\relax \let\smallcapsspacing\relax}{}{}
\makeatother


% -- Binding-aware blank page -------------------------------
\newcommand{\blankpage}{%
  \clearpage%
  \thispagestyle{empty}%
  \null\vfill%
  \begin{center}%
    \color{MidGray}\sffamily\small\itshape(This page intentionally left blank.)%
  \end{center}%
  \vfill\clearpage%
}

%\AtBeginDocument{\def\MakeLowercase#1{#1}} % TEMP: disable dangerous case-change


\usepackage{fancyhdr}
\pagestyle{fancy}
\fancyhf{}% clear all
\fancyhead[LE]{\small\scshape\leftmark}  % no MakeLowercase here
\fancyhead[RO]{\small\scshape\rightmark} % no MakeLowercase here
\fancyfoot[C]{\thepage}
\AtBeginDocument{\pagestyle{fancy}}

\hypersetup{
  colorlinks=true,
  linkcolor=blue,
  urlcolor=blue,
  citecolor=blue,
  hypertexnames=false
}

% ============================================================
%  DOCUMENT
% ============================================================
\begin{document}
\pagenumbering{alph} % Use letters for front matter (a, b, c...)
% -- Cover page ---------------------------------------------
\begin{titlepage}
    \vspace*{1.6in}
    {\sffamily\fontsize{9}{11}\selectfont\MakeUppercase{Training Manual}}
    \vspace{0.3in}
    {\color{BrandAccent}\rule{4.5in}{2pt}}
    \vspace{0.25in}
    % ADD YOUR TITLE TEXT HERE:
    {\sffamily\bfseries\fontsize{34}{40}\selectfont Switch Access \\ Training Protocol}
    \vspace{0.2in}
    {\color{BrandAccent}\rule{4.5in}{0.6pt}}
    \vspace{0.5in}
    \vfill
    {\sffamily\fontsize{7.5}{9}\selectfont\color{MidGray}%
      CONFIDENTIAL --- For Internal Use Only%
    }
  \clearpage
  \nopagecolor
\end{titlepage}

% Blank verso after cover (binding)
\blankpage

% -- Front matter -------------------------------------------
\frontmatter
\pagestyle{fancy}

{
\hypersetup{linkcolor=BrandBlue}
\tableofcontents
\clearpage
}

% -- Main content -------------------------------------------
\mainmatter
\pagenumbering{arabic} % Switch to numbers (1, 2, 3...) for the main content
\chapter{Switch access teacher training manual}
\label{switch-access-teacher-training-manual}
\chaptermark{Switch access teacher training manual}

\chapter{Teaching philosophy}
\label{section-1-introduction-teaching-philosophy}
\chaptermark{Teaching philosophy}

\section{What This Manual Is For}\label{what-this-manual-is-for}

\marginnote{This manual covers Evermusic, VLC, NLS BARD, Bookshare, Sora, and YouTube, see Sections~5, 12, and 13 for app-specific protocols.}

This manual trains educators, paraprofessionals, and support staff to
teach students to independently control music playback using
switch-based step scanning. While the immediate goal is music selection,
the real skill being built is cause-and-effect understanding,
decision-making, and environmental control, foundational capacities
for communication and independence.

This is not simply a technology tutorial. It is a teaching protocol
grounded in behavioral principles: graduated prompting, errorless
learning, data-driven decision-making, and systematic prompt fading.
Teachers who understand the \emph{why} behind each step will be better
equipped to adapt the protocol for individual students.

\section{How the Skill Is
Structured}\label{how-the-skill-is-structured}

\marginnote{Each behavior in this chain can be targeted individually. Students with severe motor fatigue may need the chain broken into even smaller steps, consult OT/PT.}

Switch scanning for music involves a chain of behaviors that students
learn progressively. The student must:

\begin{itemize}

\item
  Watch a list of songs on the screen
\item
  Press Switch 1 (Move) to advance a highlight box through the list
\item
  Recognize when the highlight is on a desired song
\item
  Press Switch 2 (Select/Pick) to start playing that song
\item
  Wait for the song to play, then independently reinitiate the cycle
\end{itemize}

Each link in this chain requires its own instruction. Students do not
typically acquire the full chain at once. This protocol teaches it piece
by piece, with the teacher providing the parts the student cannot yet do
independently and fading that support systematically.

\section{Who This Protocol Is
For}\label{who-this-protocol-is-for}

\begin{longtblr}{colspec={X[1,l] X[1,l]},rowhead=1,hlines}
\toprule
Student Profile & Appropriate Starting Stage \\
No prior switch experience & Section~3, Cause and Effect (one switch) \\
Can press a switch intentionally but has not used two switches & Section~3 \textgreater{} Phase~1 of main protocol \\
Has some two-switch experience, needs prompting & Phase~2 or~3 of main protocol \\
Independent with two switches; long lists are frustrating & Three-Switch Extension (Section~8) \\
Solid two-switch user; goal is generalization & Phase~4 + Generalization Planning \\
\bottomrule
\end{longtblr}

\section{Core Teaching
Principles}\label{core-teaching-principles}

\marginnote{These four principles apply equally to music, reading apps, and YouTube navigation. See Sections~12 and~13.}

Before beginning instruction, all teachers should internalize these
principles.


\textbf{PRINCIPLE 1: DATA DRIVES DECISIONS}

Phase advancement, prompt fading, and session length are all determined
by data, not by how many sessions have passed or how the session
felt. Record data every session. Review it before every phase decision.



\textbf{PRINCIPLE 2: PROMPT FADING IS THE GOAL}

Every prompt you provide is a temporary support on the path to
independence. The goal is always to use the least intrusive prompt that
produces a correct response, and then to fade even that. More
prompting is not more helpful; it delays independence.



\textbf{PRINCIPLE 3: MOTIVATION IS NOT OPTIONAL}

This skill is meaningless if the student does not care about the
outcome. Invest time in identifying genuinely preferred songs. A student
who is indifferent to the music content will not generalize this skill.
Consult families, observe the student, and run preference assessments.



\textbf{PRINCIPLE 4: THE MUSIC IS THE MOTIVATOR, INDEPENDENCE IS THE
SKILL}

Teachers sometimes focus so much on the music activity that they forget
they are teaching a transfer skill: the ability to control one's
environment through intentional action. Keep this in mind when
troubleshooting. If a student is ``just getting music,'' the fading is
not working.


\section{Roles on the Instructional
Team}\label{roles-on-the-instructional-team}

\marginnote{For students with visual impairments, the TVI should be consulted \emph{before} initial setup, not only when problems arise. Screen distance, contrast, and VoiceOver integration are access issues, not afterthoughts.}

\begin{longtblr}{colspec={X[1,l] X[2,l]},rowhead=1,hlines}
\toprule
Role & Responsibilities \\
Classroom Teacher & Plans sessions, tracks data, makes phase advancement decisions, communicates with families, consults specialists \\
Paraprofessional / Para & Delivers sessions following teacher direction, implements prompt levels as assigned, records data accurately, maintains device setup \\
AT Specialist & Configures switch hardware and iPadOS settings, troubleshoots technical issues, recommends switch alternatives, consults on positioning \\
OT/PT & Determines optimal switch positioning for students with motor impairments, advises on hold duration settings, manages fatigue \\
TVI (Teacher of Visually Impaired) & Advises on screen distance, contrast, cursor settings, and auditory supports for students with visual impairments \\
BCBA / Behavior Specialist & Consulted when engagement, refusal, or distress problems arise across multiple sessions \\
Families & Run home practice sessions, report observations, share preferred content, keep devices charged and connected \\
\bottomrule
\end{longtblr}

\chapter{Equipment configuration}
\label{section-2-equipment-configuration-pre-session-setup}
\chaptermark{Equipment configuration}

\section{Required Equipment}\label{required-equipment}

\marginnote{For students using NLS BARD, Bookshare, or Sora, the same switch hardware is used, no additional equipment is needed. See Section~12 for those app configurations.}
\marginnote{CAUTION: Always confirm with your district MDM administrator which apps are approved before purchasing or downloading. Some apps in this manual may require IT whitelist approval.}

\begin{longtblr}{colspec={X[1,l] X[2,l]},rowhead=1,hlines}
\toprule
Item & Notes / Specifications \\
iPad & iPadOS 16 or later; any supported iPad model \\
Bluetooth Switch Interface & e.g., Blue2, Tecla Shield, BJOY Stick, RJ Cooper Bluetooth. Must support at least 2 ports (3 for three-switch extension) \\
Switch 1, Move to Next Item & Physical switch connected to interface port 1. Color-tape or label for identification \\
Switch 2, Select / Pick & Physical switch connected to interface port 2. Different color or shape from Switch~1 \\
Evermusic App & Primary app (free). Download from App Store before first session \\
VLC for Mobile & Backup app if Evermusic is blocked by district MDM. Published by VideoLAN (open-source) \\
MP3 Files & 5-10 songs to start. Pre-loaded to local (on-device) storage only. No cloud services \\
Data Collection Sheets & Printed or digital. One sheet per session. See Section~6 for templates \\
\bottomrule
\end{longtblr}

\section{Loading Music Files onto the
iPad}\label{loading-music-files-onto-the-ipad}

\marginnote{For NLS BARD, Bookshare, and Sora (Section~12), books and audio are downloaded directly within the app, no USB transfer needed. This is one advantage of those apps over local MP3 workflows.}

All music files must be stored locally on the device. Cloud services
require login steps and internet connectivity that add unnecessary
barriers for switch users.

\textbf{Method A, USB via Finder (Mac) or iTunes (Windows)}

\begin{enumerate}
\item
  Connect iPad to computer via USB cable.
\item
  Open Finder (Mac) or iTunes (Windows).
\item
  Select your iPad, then click File Sharing.
\item
  Select Evermusic (or VLC if using that app).
\item
  Drag MP3 files into the app's document area.
\item
  Open the app on iPad and confirm files appear under Local Files.
\end{enumerate}

\textbf{Method B, AirDrop or Files App}

\begin{enumerate}
\item
  AirDrop MP3 files from a Mac or another iPad directly to the iPad.
\item
  Save to Files app \textgreater{} On My iPad \textgreater{} Evermusic
  folder.
\item
  Open Evermusic \textgreater{} Local Files to confirm.
\end{enumerate}


\textbf{FILE FORMAT NOTE:} Evermusic works best with .mp3 files. It may
not display .m4a or .aac files depending on version. Start with 5-10
songs only. A short, scannable list reduces fatigue and increases the
likelihood the student finds a preferred song quickly.


\section{iPadOS Switch Control
Configuration}\label{ipados-switch-control-configuration}

\marginnote{WARNING: Switch Control settings are stored per-device, not per-user. If the iPad is shared among students, settings must be verified and reset at the start of \emph{every} session.}

\textbf{Step 1, Pair the Bluetooth Switch Interface}

\begin{enumerate}
\item
  Go to Settings \textgreater{} Bluetooth.
\item
  Power on the switch interface. It should appear in the device list.
\item
  Tap the interface name and confirm pairing. A confirmation tone or LED
  signals success.
\item
  Activate each physical switch to confirm the iPad registers the input.
\end{enumerate}

\textbf{Step 2, Configure Switch Control}

\begin{enumerate}
\item
  Go to Settings \textgreater{} Accessibility \textgreater{} Switch
  Control.
\item
  Tap Switches \textgreater{} Add New Switch \textgreater{} External
  (Bluetooth).
\item
  Activate Switch 1 when prompted. Name it \textbf{``Next Item''} and
  set action to \textbf{Move to Next Item}. Save.
\item
  Tap Add New Switch again. Activate Switch 2 when prompted. Name it
  \textbf{``Select''} and set action to \textbf{Select Item}. Save.
\item
  Configure remaining settings per the table below.
\end{enumerate}

\begin{longtblr}{colspec={X[1,l] X[1,l] X[2,l]},rowhead=1,hlines}
\toprule
Setting & Recommended Value & Reason \\
Scanning Style & Step Scanning & Manual control, switch advances highlight; no auto-scan timer \\
Auto Scanning & OFF & Student controls pace entirely; no time pressure \\
Cursor Color & High-visibility (yellow or red) & Clear visual feedback for student \\
Cursor Size & Large & Easier to track on screen \\
Hold Duration & 0.5 seconds (starting point) & Reduces accidental activations; adjust per student motor ability \\
Ignore Repeat & ON & Prevents switch bounce triggering multiple actions from one press \\
Resume Scanning & Automatic & Highlight restarts from top of list after each selection \\
\bottomrule
\end{longtblr}

\marginnote{\textbf{Visual Impairment:} For students with low vision, set Cursor Color to bright yellow and Cursor Size to Large. Consult TVI for optimal screen distance, typically 12-18 inches for students with cortical visual impairment.}
\marginnote{Section~9 for full disability-specific adaptations including CVI, cortical visual impairment, and combined VI/motor profiles.}

\textbf{Step 3, Testing Switch Control Before Student Use}

\begin{enumerate}
\item
  Enable Switch Control (green toggle at top of Switch Control
  settings).
\item
  A colored highlight box should appear on screen.
\item
  Press Switch 1 several times, confirm highlight moves through
  items.
\item
  Press Switch 2, confirm the highlighted item activates.
\item
  If switches are reversed: return to Switches list and re-assign
  actions.
\item
  If no highlight appears: ensure Bluetooth interface is connected;
  toggle Switch Control off and on.
\end{enumerate}

\section{Evermusic Configuration (One-Time Teacher
Task)}\label{evermusic-configuration-one-time-teacher-task}

\begin{enumerate}
\item
  Open Evermusic. Tap the gear/Settings icon (top-right corner).
\item
  Under \textbf{Player}: set \textbf{Repeat Mode to Repeat Stop} (also
  shown as ``None'' or ``Stop After Song''). This is the critical
  setting, it stops playback after one song and returns to the list.
\item
  Set \textbf{Shuffle to OFF} for predictable top-to-bottom song order.
\item
  Optional: tap Main Screen Actions and hide unnecessary controls
  (Chromecast, lyrics, EQ, sleep timer). Fewer on-screen buttons = fewer
  items to scan through.
\item
  Navigate to Local Files \textgreater{} Files on this iPad
  \textgreater{} your MP3 folder. Leave this as the home screen.
\end{enumerate}

\marginnote{Hide unnecessary controls in Evermusic's Main Screen Actions. Fewer on-screen items means fewer switch presses to reach what the student needs, directly reducing fatigue for students with CP.}
\marginnote{Section~5 for VLC configuration. Section~12 for NLS BARD, Bookshare, and Sora setup.}

\textbf{CRITICAL, VERIFY REPEAT STOP IS ACTIVE BEFORE EVERY SESSION}

Evermusic may reset this setting if force-closed. Teacher: open
Evermusic \textgreater{} Settings \textgreater{} Player \textgreater{}
confirm Repeat Mode shows Stop/None before handing the iPad to the
student. If this setting is wrong, the app will auto-advance to the next
track after a song ends, bypassing the student's scanning and selection; defeating the purpose of the training.


\section{Before Every Session, Teacher Setup
Checklist}\label{before-every-session-teacher-setup-checklist}

Complete this checklist at the start of every session without exception.

\begin{itemize}

\item[\checkbox]
Bluetooth switch interface is powered on and connected to iPad
\item[\checkbox] 
Switch Control is enabled in Settings \textgreater{} Accessibility
  (green toggle)
\item[\checkbox]
Switch 1 = Move to Next Item confirmed (press it, highlight should
  advance)
\item[\checkbox] 
Switch 2 = Select Item confirmed (press it, highlighted item should
  activate)
\item[\checkbox] 
Evermusic is open and showing Local Files song list (not cloud/login
  view)
\item[\checkbox] 
Repeat Mode in Evermusic is set to Stop/None (not Repeat All or Repeat
  One)
\item[\checkbox] 
Shuffle is \textbf{OFF} in Evermusic
\item[\checkbox] 
Song list contains appropriate number of songs for student's current
  phase
\item[\checkbox] 
Data collection sheet is ready
\item[\checkbox] 
Auto-Lock set to 5 minutes or Never during session (Settings
  \textgreater{} Display and Brightness)
\end{itemize}

\chapter{One switch}
\label{section-3-prerequisite-stage-cause-effect-with-one-switch}
\chaptermark{One switch}

\marginnote{This cause-effect foundation applies to \emph{all} app contexts in this manual, not just Evermusic. Students who have not yet mastered cause-effect will struggle with BARD, Sora, and YouTube equally.}


\textbf{START HERE for students with no prior switch experience.}

This section covers the foundational stage that comes \textbf{before}
the two-switch protocol. Do not skip it for students who are new to
switches or who lack consistent cause-and-effect understanding.


\section{Purpose of the Cause-Effect
Stage}\label{purpose-of-the-cause-effect-stage}

Before a student can use switches purposefully, they must understand
that their own action, pressing the switch, directly causes
something to happen. This sounds simple, but for students who have had
limited experience controlling their environment, it is a foundational
cognitive and motor concept that must be built deliberately.

\textbf{The cause-and-effect chain being built:}


\textbf{CAUSE:} Student presses the Select switch.\\
\textbf{EFFECT:} A song the student wanted begins playing.\\
\textbf{LEARNING:} My action made something happen. I have control.


\textbf{Why navigation stays with the teacher at this stage:} The
teacher or peer helper presses Move to Next Item (or manually taps the
screen) to advance the highlight through the song list, pausing on each
song for 3-5 seconds. The student's only job is to press their switch
when they want the highlighted song.

This removes the two-task demand, watch the list AND decide when to
press a different switch to move it, and replaces it with one task:
\emph{press your switch when you want this.}

\begin{longtblr}{colspec={X[1,l] X[2,l]},rowhead=1,hlines}
\toprule
Student Profile & Why Single-Switch Works \\
Students with intellectual disabilities & Closely mirrors familiar cause-effect activities (BIGmack, single-switch toys). One demand instead of two \\
Students with visual impairments & Teacher narrates each song title aloud; student responds to auditory cue rather than tracking a visual highlight \\
Students with cerebral palsy & Isolates the Select action, student focuses all motor effort on one reliable switch movement without timing a second switch \\
\bottomrule
\end{longtblr}

\section{Setup for the Cause-Effect
Stage}\label{setup-for-the-cause-effect-stage}

\textbf{Teacher navigation options:}


\textbf{OPTION A}, Teacher uses a second physical switch assigned to
Move to Next Item. This is the more authentic setup and models switch
use for the student.

\textbf{OPTION B}, Teacher taps the song list on the screen manually,
pausing on each song for the student's response window. No additional
hardware needed.


\textbf{Song list setup:}
\begin{itemize}

\item Start with 3-5 songs maximum
\item Include at least 2 songs confirmed as preferred (use family input or preference assessments)
\item Place preferred songs in different positions, not always first, so the student must genuinely wait and respond
\item Keep song titles short and distinct for easy verbal narration
\end{itemize}

\section{Training Phases, Cause-Effect
Stage}\label{training-phases-cause-effect-stage}

\subsection{PHASE 0, Building the Cause-Effect
Connection}\label{phase-0-building-the-cause-effect-connection}

\textbf{Goal:} Student demonstrates any intentional switch press
connected to a desired outcome, even if inconsistently.

\textbf{Session Procedure:}

\begin{enumerate}
\item
  Open the song list. Say: \emph{``We're going to listen to some music.
  Press your switch when I show you a song you want.''}
\item
  Begin advancing the highlight slowly (teacher presses Move or taps
  screen). Pause on each song for \textbf{5 full seconds}.
\item
  As you pause on each song: say the song title aloud. For example:
  \emph{``\,`Happy' by Pharrell.''}
\item
  Watch the student. If they press their switch at any point during the
  5-second window: immediately select that song and let it play. Say:
  \emph{``You picked it!''}
\item
  If the student does not press during the 5-second window: move on to
  the next song \textbf{without comment}. Do not say ``no'' or redirect.
  Just move on.
\item
  After cycling through the full list once with no selection: choose a
  known preferred song yourself and play it. This ensures the session
  ends positively.
\end{enumerate}


\textbf{WHAT COUNTS AS AN INTENTIONAL PRESS IN PHASE 0}

Any press that occurs during the pause window, even if the student
does not yet clearly understand the connection. You are building the
association through repetition.

Do NOT count presses that happen while the highlight is moving (between
songs). Only count presses during the pause.

Do NOT physically prompt the switch press in Phase 0. If the student
does not press, that is data. Physical prompting here obscures whether
the student understands the cause-effect relationship.


\textbf{Phase 0 mastery criterion:} Student presses their switch during
the pause window on at least \textbf{3 out of 5} pause opportunities,
across \textbf{2 consecutive sessions}, without physical prompting.

\subsection{PHASE 1, Consistent Cause-Effect with Preferred
Choices}\label{phase-1-consistent-cause-effect-with-preferred-choices}

\textbf{Goal:} Student reliably presses their switch specifically when a
preferred song is highlighted, demonstrating true intentional
choice-making, not random pressing.

\textbf{Session Procedure:}

\begin{enumerate}
\item
  Use a list of 4-5 songs: 2-3 known preferred, 1-2 known
  non-preferred or neutral. Mix up the order each session.
\item
  Advance the highlight at the same pace as Phase 0. Pause 5 seconds on
  each song. Narrate the title.
\item
  Record whether the student presses during each pause and whether that
  song was preferred or non-preferred.
\item
  When the student presses on a preferred song: select and play it.
  Provide natural reinforcement: \emph{``You wanted that one!''}
\item
  When the student presses on a non-preferred song: \textbf{still play
  it}. Do not withhold the consequence. The learning comes from
  experiencing what their press caused.
\item
  After the song ends, return to the list and continue.
\end{enumerate}

\textbf{Phase 1 mastery criterion:} Student presses switch on \textbf{4
of 5} preferred-song pauses AND presses on \textbf{1 or fewer}
non-preferred pauses per session, across \textbf{3 consecutive
sessions}.

\subsection{PHASE 2, Reducing Wait Time and Fading Teacher
Narration}\label{phase-2-reducing-wait-time-and-fading-teacher-narration}

\textbf{Goal:} Student responds more quickly and begins to rely less on
verbal narration of song titles, preparing for visual-highlight
tracking in independent scanning.

\textbf{Session Procedure:}

\begin{enumerate}
\item
  Reduce the pause on each song from 5 seconds to \textbf{3 seconds}.
\item
  Reduce verbal narration: instead of the full title, say only the first
  word (e.g., \emph{``Happy\ldots{}''}) or fade to silence entirely over
  2-3 sessions.
\item
  If the student misses a preferred song because the window was too
  short: note it in data, but do not extend the time back. Flag for the
  IEP team if needed.
\item
  Begin occasionally using shorter windows (2-2.5 seconds) on
  non-preferred and slightly longer (3-3.5 seconds) on preferred songs
  to naturally differentiate the wait.
\end{enumerate}

\textbf{Phase 2 mastery criterion:} Student selects preferred songs with
\textbf{3-second pause and no verbal narration} on \textbf{4 of 5}
opportunities across \textbf{3 consecutive sessions}.


\textbf{PHASE 2 IS THE BRIDGE TO THE TWO-SWITCH PROTOCOL}

Once a student reliably selects preferred items during a 3-second pause
without verbal prompting, they have demonstrated: (1) intentional switch
use, (2) cause-and-effect understanding, and (3) visual attention to the
highlighted item. These are exactly the skills needed for two-switch
scanning. The transition to Section 4 can now begin.


\section{Prompting Guidelines for the Cause-Effect
Stage}\label{prompting-guidelines-for-the-cause-effect-stage}

\marginnote{WARNING: Physical prompting during cause-effect stage masks whether the student understands the relationship. Use it only for initial motor shaping in a separate, non-choice context, never during a live choice trial.}
\marginnote{Section~4.2 for the full two-switch prompting hierarchy.}


Physical prompting is \textbf{NOT recommended} in this stage. Use the
hierarchy below carefully.


\begin{longtblr}{colspec={X[1,l] X[2,l]},rowhead=1,hlines}
\toprule
Prompt Type & When and How to Use It \\
No prompt (preferred) & Default in all phases. Wait the full pause window. Record the result \\
Expectant look / lean forward & Subtle non-verbal cue. Teacher shifts posture slightly toward student as highlight lands on a preferred song. Use in Phase~0 only \\
Verbal: name the switch & \emph{``Your switch, do you want this one?''} Use sparingly in Phase~0 only. Fade by Phase~1 \\
Verbal: describe the outcome & \emph{``Press your switch if you want to hear this song.''} Use in Phase~0 only \\
Gestural: point to switch & Point to (do not touch) the student's switch while highlight is on a preferred song. Use only if verbal prompts have had no effect over 3+ sessions \\
Physical (NOT RECOMMENDED) & Physical prompting during cause-effect stage prevents you from knowing whether the student understands the relationship. If needed to build initial awareness, do this in a separate non-choice activity first \\
\bottomrule
\end{longtblr}


\textbf{IF THE STUDENT IS NOT PRESSING AFTER 3-5 SESSIONS WITH NO
PROMPTS}

Do not increase prompting. Instead, investigate:
\begin{itemize}

\item Are the songs actually preferred? Get fresh preference data from the family.
\item Can the student physically access the switch? Consult OT/PT.
\item Is the pause window long enough? Try extending to 7-10 seconds.
\item Is the cause-effect concept established in other contexts? Try a simpler single-switch toy (press switch $\rightarrow$ toy activates) with no choice involved.
\item Is there a sensory barrier? Can the student see the screen and hear the music adequately?
\end{itemize}


\section{Transition Readiness Checklist, Moving to Two-Switch
Protocol}\label{transition-readiness-checklist-moving-to-two-switch-protocol}

All items must be confirmed before transitioning to Section 4.

\begin{itemize}

\item[\checkbox] 
Student presses Select switch on 4 of 5 preferred-song pauses without
  physical prompting (Phase 2 criterion met across 3 consecutive
  sessions)
\item[\checkbox] 
Student shows clear discrimination: presses significantly more on
  preferred songs than non-preferred songs (discrimination index $\geq$ 60\%)
\item[\checkbox] 
Student does not require verbal narration to respond (narration fully
  faded in Phase 2)
\item[\checkbox] 
Student tolerates 3-second pause window without significant distress
  or fatigue within a standard session
\item[\checkbox] 
Switch is reliably accessible, positioning confirmed by OT/PT if
  student has motor complexity
\item[\checkbox] 
Teacher has reviewed Section 4 of this manual and is ready to begin
  Phase 1
\item[\checkbox] 
Data from this stage has been documented and shared with the IEP team
\end{itemize}

\chapter{Two-switch training}
\label{section-4-two-switch-training-protocol}
\chaptermark{Two-switch training}

\section{The Full Session
Workflow}\label{the-full-session-workflow}

\marginnote{This 7-step cycle is the \emph{same} cognitive structure used in YouTube navigation (Section~13) and BARD/Bookshare/Sora (Section~12), scan, select, engage with content, return. Generalization across apps builds on this same chain.}

Once both switches are in play, the complete interaction loop looks like
this:

\begin{longtblr}{colspec={X[1,l] X[7,l]},rowhead=1,hlines}
\toprule
Step & Action \\
\textbf{Step 1} & Song list is displayed in Evermusic (Local Files view) \\
\textbf{Step 2} & Student presses Switch~1 (Next Item), highlight moves down the list \\
\textbf{Step 3} & Student continues pressing Switch~1 until highlight is on the desired song \\
\textbf{Step 4} & Student presses Switch~2 (Select), song begins playing \\
\textbf{Step 5} & Song plays to completion (no auto-advance, Repeat Stop is active) \\
\textbf{Step 6} & Evermusic returns to song list automatically \\
\textbf{Step 7} & Student repeats from Step~2 to choose the next song, no prompt needed \\
\bottomrule
\end{longtblr}


\textbf{KEY POINT: NO AUTO-PLAY}

After each song ends, the student must actively scan and select again.
This is intentional, it provides repeated practice opportunities for
switch scanning within a motivating activity. If a student finds this
frustrating, reduce the number of songs in the list so they reach
preferred songs more quickly.


\section{Prompting Hierarchy}\label{prompting-hierarchy}

\marginnote{Always begin each session at the student's \emph{last successful prompt level} from the previous session, not the highest level you know works. Consistency is more important than speed of fading.}
\marginnote{CAUTION: For students with high spasticity (CP), ensure switch placement is confirmed by OT/PT before advancing prompt levels. Motor access variability can look like ``inconsistency'' when it is actually a positioning problem.}

\begin{longtblr}{colspec={X[1,l] X[2,l]},rowhead=1,hlines}
\toprule
Prompt Level & How to Apply It \\
Full physical guidance (Phase~1) & Put your hand over the student's hand and press the switch together. Say the action out loud: \emph{``Move to next.'' / ``Pick that song.''} \\
Partial physical, wrist/hand-under-hand & Rest hand lightly on student's wrist as a cue, do not move their hand. Point to the switch and wait 3 seconds \\
Gestural (Phase~2) & Point to the correct switch without touching. Wait 3 seconds \\
Verbal only (Phase~3) & Say: \emph{``Find a song.''} Wait. Only speak again if 10 seconds pass with no movement \\
Independent (Phase~4) & Say: \emph{``Music time.''} Then stay quiet and let the student drive completely \\
\bottomrule
\end{longtblr}

\section{Phase 1, Errorless Learning / Full Physical
Guidance}\label{phase-1-errorless-learning-full-physical-guidance}

\textbf{Goal:} Student experiences the complete switch action and
response without errors. Teacher physically guides hand to each switch.

\textbf{Session Setup:}
\begin{itemize}

\item Arrange 3-5 songs in the Evermusic list (very short list)
\item Sit beside student, on the same side as dominant hand
\item Verbally label switches before starting: \emph{``This switch moves to the next song. This switch picks the song.''}
\end{itemize}

\textbf{Trial Procedure:}

\begin{enumerate}
\item
  Point to song list on screen. Say: \emph{``Let's find a song to listen
  to.''}
\item
  Physically guide student's hand to Switch 1. Activate together. Say:
  \emph{``Move to next song.''} Observe highlight advance.
\item
  Repeat Step 2 one or two times to demonstrate scanning.
\item
  Physically guide student's hand to Switch 2. Activate together. Say:
  \emph{``Pick that song.''} Song begins playing.
\item
  Allow song to play (30+ seconds recommended).
\item
  Song ends. Evermusic returns to list. Say: \emph{``The song stopped.
  Let's find another one.''}
\item
  Repeat for 3-5 songs per session.
\end{enumerate}

\textbf{Phase 1 mastery criterion:} Student tolerates physical guidance
and does not resist switch activation across \textbf{2 consecutive
sessions}.

\section{Phase 2, Partial Physical Guidance with Gestural
Prompts}\label{phase-2-partial-physical-guidance-with-gestural-prompts}

\textbf{Goal:} Reduce physical guidance. Use hand-on-wrist prompt, then
fade to pointing/gestural only.

\begin{enumerate}
\item
  Follow same session setup as Phase 1 (3-5 songs).
\item
  Instead of full hand-over-hand: rest hand lightly on student's wrist
  as a cue, do not move their hand.
\item
  Point to Switch 1 and say: \emph{``Move to next song.''} Wait 3
  seconds.
\item
  If student activates Switch 1 independently: provide immediate natural
  reinforcement (\emph{``You moved to the next song!''}).
\item
  If no response after 3 seconds: provide light wrist prompt, then back
  off on the next trial.
\item
  After student activates Switch 1 one or two times, point to Switch 2:
  \emph{``Pick that song.''} Wait 3 seconds.
\item
  Follow the same prompting hierarchy for Switch 2.
\end{enumerate}

\textbf{Phase 2 mastery criterion:} Student independently activates each
switch with gestural prompt only (point/verbal) on \textbf{4 of 5}
consecutive trials.

\section{Phase 3, Verbal Prompt
Only}\label{phase-3-verbal-prompt-only}

\textbf{Goal:} Student scans and selects with only a verbal cue, no
physical or gestural prompt.

\begin{enumerate}
\item
  Expand song list to 5-10 songs.
\item
  Present iPad. Say only: \emph{``Find a song you want.''}
\item
  Wait. Observe student scan with Switch 1.
\item
  If student pauses on a song and activates Switch 2: provide
  reinforcement.
\item
  If student scans past all songs without selecting: say \emph{``Did you
  find one you want?''}, do not point or gesture.
\item
  If student does not activate Switch 2 after 10 seconds on a song:
  provide one verbal reminder only: \emph{``Press the pick switch to
  choose that song.''}
\item
  After song ends and list reappears: do \textbf{not} prompt. Wait 10
  seconds to see if student initiates next scan independently.
\end{enumerate}

\textbf{Phase 3 mastery criterion:} Student completes full
scan-and-select cycle (find song $\rightarrow$ pick song $\rightarrow$ listen to completion $\rightarrow$
return to list and re-initiate scanning) with verbal prompt only on
\textbf{4 of 5} trials across \textbf{2 sessions}.

\section{Phase 4, Independent
Performance}\label{phase-4-independent-performance}

\textbf{Goal:} No prompts. Student drives the entire interaction from
start to finish.

\begin{enumerate}
\item
  Expand song list to 10-20 songs if appropriate for the student.
\item
  Present iPad. Say nothing except what would be natural (e.g.,
  \emph{``Here's your music time.''}).
\item
  Do not provide any prompts unless student makes a switch error
  (activating Select when they meant to scan further).
\item
  If a switch error occurs, allow the natural consequence (a song plays
  they didn't want), then observe if the student independently
  navigates back to the list.
\item
  Record data: see Section 6 for data collection templates.
\end{enumerate}

\textbf{Phase 4 mastery criterion:} Student completes \textbf{3 full
independent} scan-select-listen-return cycles per session with
\textbf{zero prompts} across \textbf{3 consecutive sessions}.

\section{Between-Song Moments, How to Handle the
Pause}\label{between-song-moments-how-to-handle-the-pause}

The pause between songs (when the list reappears after a song ends) is
one of the most instructionally important moments in the session. Do not
rush it.

\begin{longtblr}{colspec={X[1,l] X[1,l]},rowhead=1,hlines}
\toprule
What the Student Does & Teacher Response \\
Immediately begins scanning again & Say nothing. This is the goal. Record ``independent re-initiation.'' \\
Pauses 5-10 seconds, then begins scanning & Say nothing. Still counts as independent. \\
Does not begin after 10+ seconds & Remind once: \emph{``Find another song if you want.''} Then wait again. \\
Looks away or turns from iPad & That is okay. Music time may be over. Do not push for another choice. \\
Shows distress & End the session immediately. See Section~7 for behavioral troubleshooting. \\
\bottomrule
\end{longtblr}


\textbf{WHAT COUNTS AS SUCCESS}

\begin{itemize}

\item
  Student pressed a switch independently, even once. That is
  progress.
\item
  Student tolerated the session without distress. That matters.
\item
  Student showed any sign of preference, leaning toward a song,
  looking at the screen, a smile when music played.
\end{itemize}

You do NOT need the student to complete the full cycle independently
right away. Partial steps are wins.


\section{Generalization Goals}\label{generalization-goals}

\marginnote{App generalization is a natural generalization target: a student who scans a music list in Evermusic is ready to scan a playlist in YouTube (Section~13) or a book list in BARD (Section~12). Introduce new app contexts as generalization goals in the IEP.}
\marginnote{Sections~12 and~13 for reading app and YouTube generalization targets. Section~11.2 for IEP integration notes.}

Once mastery is achieved in the primary training setting, systematically
expand to new contexts. \textbf{Change only one variable at a time:}

\begin{itemize}

\item
  Different song lists (different folder of MP3s)
\item
  Different time of day (not just dedicated music time)
\item
  Different communication partners present (substitute teacher, peer,
  family member)
\item
  Different physical environments (classroom, library, home, community)
\item
  Student independently requesting music time using AAC device or
  symbol, then completing the full scanning sequence
\end{itemize}

\section{\texorpdfstring{\NoCaseChange{VLC alternate configuration}}{VLC alternate configuration}}
\label{section-5-vlc-alternate-configuration}

\marginnote{NLS BARD, Bookshare, and Sora (Section~12) are also typically approved on district MDM profiles, often more readily than Evermusic, because they are federally authorized or nonprofit education-focused apps.}


\textbf{USE THIS SECTION WHEN EVERMUSIC IS BLOCKED BY YOUR DISTRICT'S
MDM}

VLC for iOS/iPadOS is open-source, free, and widely approved on district
MDM profiles due to its nonprofit/educational reputation. It stores and
plays files entirely locally, no account, no cloud, no login
required.


\section{Installing VLC and Loading
Files}\label{installing-vlc-and-loading-files}

\begin{enumerate}
\item
  Search the App Store for \textbf{``VLC for Mobile''} (developer:
  VideoLAN). Install.
\item
  Load MP3 files using the same USB/Finder method from Section 2.2 --
  select VLC under File Sharing and drag in MP3 files.
\item
  Alternatively: save files to Files app \textgreater{} On My iPad
  \textgreater{} VLC folder directly.
\item
  Confirm files are visible: Open VLC \textgreater{} tap the Files icon
  (folder icon, bottom nav bar). MP3s should appear as a vertical list.
\end{enumerate}

\section{Key Differences from
Evermusic}\label{key-differences-from-evermusic}

\begin{longtblr}{colspec={X[1,l] X[1,l] X[1,l]},rowhead=1,hlines}
\toprule
Feature & Evermusic & VLC for iPadOS \\
Cost & Free (basic version) & Free, open-source \\
District MDM approval & Varies, may be blocked & Usually approved (open-source nonprofit) \\
Local files only (no cloud) & Yes & Yes \\
Auto-return to list after song & Yes (with Repeat Stop) & No, student must navigate back \\
Extra steps per song cycle & 0 (list auto-reappears) & 2-3 extra switch presses (back to list) \\
Repeat Stop setting & Explicit setting in Player menu & Toggle on now-playing screen \\
Training complexity & Lower & Slightly higher (return-to-list step) \\
Best for & Primary choice; simpler workflow & When Evermusic is MDM-blocked \\
\bottomrule
\end{longtblr}

\section{VLC Session Workflow (8-Step
Cycle)}\label{vlc-session-workflow-8-step-cycle}

\marginnote{CAUTION: The return-to-list step (Steps~6-7) adds 2-3 extra switch presses per song cycle. For students with significant motor fatigue, this overhead matters. Prefer Evermusic when both are available.}
\marginnote{Section~5.5 for isolated teaching of the return-to-list step.}

\begin{longtblr}{colspec={X[1,l] X[7,l]},rowhead=1,hlines}
\toprule
Step & Action \\
\textbf{Step 1} & VLC is open on the Files tab, song list visible as a vertical list \\
\textbf{Step 2} & Student presses Switch~1 (Next Item), highlight moves down the file list \\
\textbf{Step 3} & Student presses Switch~1 until highlight is on desired song \\
\textbf{Step 4} & Student presses Switch~2 (Select), VLC opens now-playing screen and song begins \\
\textbf{Step 5} & Song plays to completion. Now-playing screen remains visible (progress bar at end) \\
\textbf{Step 6} & Student presses Switch~1 to scan to the back/list button (arrow icon, top-left of player) \\
\textbf{Step 7} & Student presses Switch~2 (Select) to return to the file list \\
\textbf{Step 8} & File list reappears. Student repeats from Step~2 \\
\bottomrule
\end{longtblr}


\textbf{CRITICAL, LEAVE VLC ON THE FILES TAB BEFORE HANDING TO
STUDENT}

VLC remembers the last screen. If it was left on the now-playing screen
from the previous session, Switch Control will scan the player controls
rather than the song list.

\textbf{Teacher action:} Before each session, open VLC and tap the Files
icon (bottom nav bar) to confirm the song list is visible. Then hand the
iPad to the student.


\section{VLC Pre-Session
Checklist}\label{vlc-pre-session-checklist}

\begin{itemize}

\item[\checkbox] 
Bluetooth switch interface is powered on and connected to iPad
\item[\checkbox] 
Switch Control is enabled in Settings \textgreater{} Accessibility
\item[\checkbox] 
Switch 1 = Move to Next Item confirmed (test: highlight advances)
\item[\checkbox] 
Switch 2 = Select Item confirmed (test: highlighted item activates)
\item[\checkbox] 
VLC is open and showing the Files tab / song list (not the now-playing
  screen)
\item[\checkbox] 
Repeat is OFF in VLC (verify on now-playing screen: repeat icon not
  highlighted)
\item[\checkbox] 
Shuffle is OFF in VLC
\item[\checkbox] 
Song list contains appropriate number of songs for student's current
  training phase
\item[\checkbox] 
Student has been taught the return-to-list step (Section 5.5) before
  independent trials
\item[\checkbox] 
Data collection sheet is ready
\end{itemize}

\section{Teaching the Return-to-List Step (VLC
Only)}\label{teaching-the-return-to-list-step-vlc-only}

Because VLC requires an active navigation step back to the file list,
this skill should be taught as a discrete component \textbf{before} full
independent trials begin.

\textbf{Isolated Practice Procedure:}

\begin{enumerate}
\item
  Teacher plays a short clip (selects a song manually). Pauses it.
\item
  Present the now-playing screen to the student.
\item
  Prompt: \emph{``The song is done. Find the list.''} Use the same
  prompting hierarchy from Section 4.
\item
  Target: student independently scans to the back/list arrow and selects
  it within 15 seconds.
\item
  \textbf{Criterion before full independent trials:} 4 of 5 correct
  returns with verbal prompt only.
\end{enumerate}

The back/list arrow is located at the \textbf{top-left} of the
now-playing screen. In Switch Control step scanning, it is typically one
of the first items the highlight lands on when scanning begins from the
top of the screen. Verify this on your specific device and VLC version
before training.

\chapter{Data collection}
\label{section-6-data-collection}
\chaptermark{Data collection}

\section{What to Measure in Two-Switch
Sessions}\label{what-to-measure-in-two-switch-sessions}

\begin{longtblr}{colspec={X[1,l] X[1,l]},rowhead=1,hlines}
\toprule
Data Point & How to Record \\
Switch~1 activations to find a song & Tally count per song selection (lower = more efficient scanning) \\
Prompt level for Switch~1 & I = Independent, V = Verbal, G = Gestural, P = Physical \\
Prompt level for Switch~2 & I = Independent, V = Verbal, G = Gestural, P = Physical \\
Correct switch use (no errors) & Y / N per trial (N = activated wrong switch) \\
Independent re-initiation after song end & Y / N, did student start scanning again without prompt? \\
Songs selected per session & Count (indicates engagement and motivation) \\
\bottomrule
\end{longtblr}

\clearpage
\section{DATA SHEET A, Two-Switch Session Data
Sheet}\label{data-sheet-a-two-switch-session-data-sheet}

\emph{(Reproducible, print one per session)}

\textbf{Student:} \_\_\_\_\_\_\_\_\_\_\_\_\_\_\_\_\_\_\_\_\_\_\_\_\_
\textbf{Date:} \_\_\_\_\_\_\_\_\_\_\_\_ \textbf{Session \#:} \_\_\_\_
\textbf{Phase:} \_\_\_\_

\textbf{Teacher/Para:} \_\_\_\_\_\_\_\_\_\_\_\_\_\_\_\_\_\_\_
\textbf{Song List Version:}
\_\_\_\_\_\_\_\_\_\_\_\_\_\_\_\_\_\_\_\_\_\_\_\_\_

\begin{longtblr}{colspec={X[1,c] X[3,l] X[3,l] X[3,l] X[3,l] X[3,l] X[3,l]},rowhead=1,hlines}
\toprule
Trial & SW1 Act. (tally) & SW1 Prompt (I/V/G/P) & SW2 Prompt (I/V/G/P) & Correct? (Y/N) & Re-init.? (Y/N) & Notes \\
1 & & & & & & \\
2 & & & & & & \\
3 & & & & & & \\
4 & & & & & & \\
5 & & & & & & \\
6 & & & & & & \\
7 & & & & & & \\
8 & & & & & & \\
9 & & & & & & \\
10 & & & & & & \\
11 & & & & & & \\
12 & & & & & & \\
\bottomrule
\end{longtblr}

\textbf{Session Totals:}

\begin{itemize}

\item
  Total trials: \_\_\_\_
\item
  Trials with both switches independent (I/I): \_\_\_\_
\item
  Songs selected: \_\_\_\_
\item
  Re-initiations without prompt: \_\_\_\_
\item
  Technical issues (describe):
  \_\_\_\_\_\_\_\_\_\_\_\_\_\_\_\_\_\_\_\_\_\_\_\_\_\_\_\_\_\_\_\_\_\_\_\_\_\_\_\_
\end{itemize}
\clearpage
\section{DATA SHEET B, Cause-Effect Stage Data
Sheet}\label{data-sheet-b-cause-effect-stage-data-sheet}

\emph{(Reproducible, print one per session)}

\textbf{Student:} \_\_\_\_\_\_\_\_\_\_\_\_\_\_\_\_\_\_\_\_\_\_\_\_\_
\textbf{Date:} \_\_\_\_\_\_\_\_\_\_\_\_ \textbf{Phase:} 0 / 1 / 2
\emph{(circle one)}

\begin{longtblr}{colspec={X[3,l] X[1,l] X[2,l] X[2,l] X[3,l]},rowhead=1,hlines}
\toprule
Song Title & P / NP & Pressed? (Y/N) & Timing (E/L/N) & Notes \\
 & & & & \\
 & & & & \\
 & & & & \\
 & & & & \\
 & & & & \\
 & & & & \\
 & & & & \\
 & & & & \\
 & & & & \\
 & & & & \\
\bottomrule
\end{longtblr}

\textbf{Key:} P = Preferred \quad NP = Non-preferred \quad
Timing: E = Early (first 2 sec) / L = Late (after 2 sec) / N = No press

\textbf{Session Totals:}

\begin{itemize}

\item
  Preferred presented: \_\_\_\_ / Preferred pressed: \_\_\_\_
\item
  Non-preferred presented: \_\_\_\_ / Non-preferred pressed: \_\_\_\_
\item
  \textbf{Preferred hit rate:} \_\_\_\_\% (preferred pressed ÷ preferred
  presented)
\item
  \textbf{Non-preferred press rate:} \_\_\_\_\% (non-preferred pressed ÷
  non-preferred presented)
\item
  \textbf{Discrimination index:} \_\_\_\_\% \textminus \_\_\_\_\% = \_\_\_\_\%
  \emph{(target: $\geq$ 60\%)}
\end{itemize}

\emph{Example: Student pressed on 4 of 5 preferred songs (80\%) and 1 of
3 non-preferred songs (33\%). Discrimination index = 80\% \textminus 33\% = 47\%.
Not yet at criterion, continue Phase 1.}

\clearpage
\section{DATA SHEET C, Three-Switch}\label{data-sheet-c-three-switch-additional-tracking}

\emph{(Use this sheet in addition to Data Sheet A once three-switch
protocol begins)}

\textbf{Student:} \_\_\_\_\_\_\_\_\_\_\_\_\_\_\_\_\_\_\_\_\_\_\_\_\_
\textbf{Date:} \_\_\_\_\_\_\_\_\_\_\_\_ \textbf{Three-Switch Phase:}
\_\_\_\_

\begin{longtblr}{colspec={X[1,c] X[3,l] X[3,l] X[3,l] X[3,l] X[4,l]},rowhead=1,hlines}
\toprule
Trial & SW3 Used? (Y/N) & SW3 Purposeful? (Y/N) & SW3 Error? (Y/N) & SW3 Prompt (I/V/G/P) & Total Presses (all switches) \\
1 & & & & & \\
2 & & & & & \\
3 & & & & & \\
4 & & & & & \\
5 & & & & & \\
6 & & & & & \\
7 & & & & & \\
8 & & & & & \\
\bottomrule
\end{longtblr}

\textbf{Session Totals:}

\begin{itemize}

\item
  SW3 activations: \_\_\_\_ \quad Purposeful uses: \_\_\_\_ \quad Errors: \_\_\_\_
\item
  Average total presses per song selection: \_\_\_\_ \emph{(should
  decrease over time)}
\end{itemize}

\textbf{Interpretation notes:}
\begin{itemize}

\item High rate of SW3 errors $\rightarrow$ may not be ready for three switches; consider returning to two-switch system
\item Very low SW3 use is \textbf{not a problem} if other goals are met, the switch is an available tool, not a required step
\item Decreasing total press count over time = increasing scanning efficiency (the key long-term metric)
\end{itemize}

\chapter{Troubleshooting}
\label{section-7-troubleshooting-guide}
\chaptermark{Troubleshooting}

\marginnote{This troubleshooting guide covers Evermusic and VLC. For NLS BARD, Bookshare, Sora, and YouTube-specific issues, see the troubleshooting subsections within Sections~12 and~13.}

\emph{Organized by what you observe going wrong. Find the matching
symptom, work through fixes in order, and stop when the problem
resolves. marks the most commonly needed fixes, try these first.}

\section{Switch Hardware
Problems}\label{switch-hardware-problems}

\subsection{Switch presses do nothing, no response at
all}\label{switch-presses-do-nothing-no-response-at-all}

\begin{enumerate}
\item
  Confirm the Bluetooth switch interface is powered on (check LED
  indicator).
\item
  On iPad: Settings \textgreater{} Bluetooth. Is the interface listed as
  ``Connected''? If ``Not Connected'' or missing: tap it to reconnect,
  or toggle Bluetooth off and on.
\item
  Confirm Switch Control is enabled: Settings \textgreater{}
  Accessibility \textgreater{} Switch Control (green toggle).
\item
  Test each physical switch, if highlight does not move, check
  physical switch cable connections to the interface.
\item
  Battery: charge the interface if the LED is dim or flashing.
\item
  Re-pair if needed: Settings \textgreater{} Bluetooth \textgreater{}
  tap ``i'' next to device \textgreater{} Forget Device. Power cycle the
  interface, then re-pair.
\end{enumerate}

\subsection{Switch activates but the wrong action
happens}\label{switch-activates-but-the-wrong-action-happens}

\emph{Symptom: Switch 1 (Move) plays a song; Switch 2 (Pick) advances
the highlight.}

\begin{enumerate}
\item
  Settings \textgreater{} Accessibility \textgreater{} Switch Control
  \textgreater{} Switches.
\item
  Tap each switch entry and confirm: Switch 1 = Move to Next Item,
  Switch 2 = Select Item.
\item
  If reversed: re-assign. Tap the entry \textgreater{} change action
  \textgreater{} save.
\item
  Test again before returning iPad to student.
\end{enumerate}

\subsection{Switch registers multiple presses from one activation
(bounce)}\label{switch-registers-multiple-presses-from-one-activation-bounce}

\emph{Symptom: One tap moves the highlight 2-3 items at once, or plays
a song immediately after moving.}

\begin{enumerate}
\item
  Settings \textgreater{} Accessibility \textgreater{} Switch Control
  \textgreater{} Ignore Repeat. Enable this setting.
\item
  If already on: increase the Ignore Repeat interval to 0.5-1.0
  seconds.
\item
  If problem persists: try swapping in a different physical switch.
\item
  Check cable connections between switch and interface, a loose
  connection can cause intermittent contacts.
\end{enumerate}

\subsection{Switch requires very hard activation or is
inconsistent}\label{switch-requires-very-hard-activation-or-is-inconsistent}

\begin{enumerate}
\item
  Check that the physical switch cable is fully seated in the interface
  jack.
\item
  Test with a different switch cable.
\item
  Check whether the switch type is appropriate for the student's motor
  ability, consult OT/PT if the student has CP or significant motor
  impairment.
\item
  Reduce Hold Duration: Settings \textgreater{} Accessibility
  \textgreater{} Switch Control \textgreater{} Hold Duration. Lower
  value = easier to trigger.
\end{enumerate}

\section{Switch Control (iPadOS)
Problems}\label{switch-control-ipados-problems}

\subsection{The highlight box does not appear on
screen}\label{the-highlight-box-does-not-appear-on-screen}

\begin{enumerate}
\item
  Confirm Switch Control is ON: Settings \textgreater{} Accessibility
  \textgreater{} Switch Control (must be green).
\item
  If ON but no highlight: press Switch 1 once. Sometimes the overlay
  needs a trigger to appear after app changes.
\item
  If highlight is present but invisible: Settings \textgreater{}
  Accessibility \textgreater{} Switch Control \textgreater{} Cursor
  Color. Change to yellow, red, or bright green. Increase Cursor Size to
  Large.
\item
  If the iPad was locked and just unlocked: Switch Control resumes when
  you activate a switch.
\end{enumerate}

\subsection{Highlight scans automatically without student pressing a
switch}\label{highlight-scans-automatically-without-student-pressing-a-switch}

\emph{Cause: Auto Scanning is enabled. This is the wrong mode for
two-switch step scanning.}

\begin{enumerate}
\item
  Settings \textgreater{} Accessibility \textgreater{} Switch Control
  \textgreater{} Auto Scanning: turn \textbf{OFF}.
\item
  Settings \textgreater{} Accessibility \textgreater{} Switch Control
  \textgreater{} Scanning Style: set to \textbf{Step Scanning}.
\item
  Test: highlight should only move when Switch 1 is pressed.
\end{enumerate}

\subsection{Switch Control turns itself off
mid-session}\label{switch-control-turns-itself-off-mid-session}

\begin{enumerate}
\item
  Add Switch Control to Accessibility Shortcut for quick re-enabling:
  Settings \textgreater{} Accessibility \textgreater{} Accessibility
  Shortcut \textgreater{} Switch Control.
\item
  To re-enable quickly: \textbf{triple-click} the Home or Side button.
\item
  Prevent lock during session: Settings \textgreater{} Display and
  Brightness \textgreater{} Auto-Lock \textgreater{} set to 5 minutes or
  Never during sessions.
\end{enumerate}

\subsection{Too many items are scanned on the now-playing
screen}\label{too-many-items-are-scanned-on-the-now-playing-screen}

\emph{Symptom: After selecting a song, the highlight scans through many
small buttons, making navigation confusing.}

\begin{enumerate}
\item
  In Evermusic: Settings \textgreater{} Main Screen Actions. Hide
  unnecessary controls (Chromecast, lyrics, EQ, sleep timer). Keep only
  Play/Pause.
\item
  Consider whether the student actually needs to interact with the
  player screen during early phases, if the song plays to completion
  automatically, they may not need to touch anything until the list
  reappears.
\end{enumerate}

\section{Evermusic Problems}\label{evermusic-problems}

\subsection{[Priority] Song auto-advances to next track (does not stop after
one
song)}\label{song-auto-advances-to-next-track-does-not-stop-after-one-song}

\emph{Cause: Repeat Mode is not set to Stop.}

\begin{enumerate}
\item
  Open Evermusic \textgreater{} tap gear icon (Settings) \textgreater{}
  Player.
\item
  Set Repeat Mode to \textbf{Stop} (sometimes labeled ``None'' or
  ``Repeat Off'').
\item
  Verify: play a short test track, let it finish, confirm the list
  reappears.
\end{enumerate}

\emph{Note: Evermusic may reset this setting if force-closed. Check it
at the start of every session.}

\subsection{Song list is empty or shows ``No files
found''}\label{song-list-is-empty-or-shows-no-files-found}

\begin{enumerate}
\item
  Tap Local Files \textgreater{} Files on this iPad. Navigate to the
  folder where MP3s were saved.
\item
  If the folder is empty: re-transfer via USB/Finder: connect iPad
  \textgreater{} Finder \textgreater{} your iPad \textgreater{} File
  Sharing \textgreater{} Evermusic \textgreater{} drag MP3s in.
\item
  Make sure files end in \textbf{.mp3}. Evermusic may not display .m4a
  or .aac depending on version.
\item
  Try closing and re-opening Evermusic after adding files.
\end{enumerate}

\subsection{Evermusic shows a cloud storage or login
screen}\label{evermusic-shows-a-cloud-storage-or-login-screen}

\begin{enumerate}
\item
  Tap the house icon or \textbf{Local Files} tab in the bottom
  navigation bar.
\item
  Do not sign in to any cloud service, all files should be local.
\item
  If the app was updated and defaulted to a cloud view: navigate
  manually to Local Files and this will become the default going
  forward.
\end{enumerate}

\section{VLC-Specific Problems}\label{vlc-specific-problems}

\subsection{[Priority] VLC opens to now-playing screen instead of file
list}\label{vlc-opens-to-now-playing-screen-instead-of-file-list}

\emph{Cause: VLC remembers the last screen from the previous session.}

\begin{enumerate}
\item
  Before handing the iPad to the student: open VLC, tap the Files icon
  (folder icon in bottom nav bar).
\item
  The song list should now be visible. Make this part of the pre-session
  checklist every time.
\end{enumerate}

\subsection{Song loops in VLC instead of
stopping}\label{song-loops-in-vlc-instead-of-stopping}

\emph{Cause: Repeat was left ON from a previous session.}

\begin{enumerate}
\item
  From the now-playing screen: locate the repeat icon (two arrows). It
  should not be highlighted.
\item
  Use Switch Control to scan to the repeat icon and deselect it, or tap
  it manually before handing the iPad to the student.
\end{enumerate}

\subsection{Student is stuck on the now-playing screen and cannot
return to
list}\label{student-is-stuck-on-the-now-playing-screen-and-cannot-return-to-list}

\begin{enumerate}
\item
  Provide verbal prompt: \emph{``Find the list button.''}
\item
  The back/list button is the \textbf{left-facing arrow at the top-left}
  of the VLC player screen.
\item
  If this happens repeatedly: run isolated return-to-list practice
  (Section 5.5) before resuming full sessions.
\end{enumerate}

\subsection{Files do not appear in VLC file
browser}\label{files-do-not-appear-in-vlc-file-browser}

\begin{enumerate}
\item
  Open the iOS Files app. Navigate to On My iPad \textgreater{} VLC.
  Confirm MP3 files are present.
\item
  If absent: transfer files. Connect iPad to computer \textgreater{}
  Finder \textgreater{} your iPad \textgreater{} File Sharing
  \textgreater{} VLC \textgreater{} drag MP3s in.
\item
  Re-open VLC \textgreater{} Files tab. Files should now appear.
\item
  If VLC does not appear under File Sharing: ensure VLC is fully
  installed and has been opened at least once.
\end{enumerate}

\section{Student Behavior and Engagement
Problems}\label{student-behavior-and-engagement-problems}


These are not technical problems, they are signals from the student
about their current access, fatigue, understanding, or motivation level.
Address them by adjusting the environment or session design, not by
increasing prompts.


\subsection{Student presses Switch 1 repeatedly without stopping to
select}\label{student-presses-switch-1-repeatedly-without-stopping-to-select}

\emph{This is usually exploration, not non-compliance.}

\begin{itemize}

\item
  Wait, allow up to 2-3 full cycles through the list.
\item
  If it continues: provide one verbal cue: \emph{``Pick one when you're
  ready.''}
\item
  Consider whether the song list is too long, reduce to 5 songs or
  fewer.
\item
  If this is consistent across sessions: consult teacher. The student
  may need a shorter list or different preferred content.
\end{itemize}

\subsection{Student shows no response to either
switch}\label{student-shows-no-response-to-either-switch}

\begin{itemize}

\item
  Check positioning first, is the switch placed where the student can
  physically reach it?
\item
  Check motivation, are the songs actually preferred? Ask family for
  home favorites.
\item
  Check sensory access, can the student see/hear the content well
  enough to be motivated by it?
\item
  Return to Phase 1 (full physical guidance) temporarily.
\item
  Shorten sessions to 2-3 minutes to reduce fatigue.
\item
  Consult OT/PT if motor access to the switch seems to be the barrier.
\end{itemize}

\subsection{Student becomes distressed during
scanning}\label{student-becomes-distressed-during-scanning}

\begin{itemize}

\item
  Stop the session immediately. Do not try to finish the trial.
\item
  Identify what triggered the distress: was it the wait time, the switch
  activation, the sound, or the environment?
\item
  Reduce the demand: shorter wait time, hand-over-hand guidance, fewer
  songs.
\item
  For students with significant anxiety or sensory sensitivities:
  consult the BCBA before continuing.
\item
  Document the trigger and response for the IEP/behavior plan.
\end{itemize}

\subsection{Student is accurate in sessions but does not generalize
to new
settings}\label{student-is-accurate-in-sessions-but-does-not-generalize-to-new-settings}

\begin{itemize}

\item
  Systematically introduce new settings one variable at a time:
  different room, different time of day, different communication
  partner.
\item
  Keep all other variables the same when introducing a new setting.
\item
  Fade to the new setting gradually over 3-5 sessions.
\item
  See Section 4.8 for generalization planning.
\end{itemize}

\subsection{Student with visual impairment cannot track the
highlight}\label{student-with-visual-impairment-cannot-track-the-highlight}

\begin{itemize}

\item
  Increase cursor size: Settings \textgreater{} Accessibility
  \textgreater{} Switch Control \textgreater{} Cursor Size
  \textgreater{} Large.
\item
  Change cursor color to maximum contrast: bright yellow on dark, or
  bright red.
\item
  Provide an auditory description of each item as the highlight lands on
  it (teacher or para narrates the song title).
\item
  Consult the TVI for screen distance and angle recommendations.
\item
  Consider whether VoiceOver integration with Switch Control could
  provide spoken item names automatically, discuss with TVI and AT
  specialist.
\end{itemize}

\section{Escalation, When to Contact Other Team
Members}\label{escalation-when-to-contact-other-team-members}

\begin{longtblr}{colspec={X[1,l] X[1,l]},rowhead=1,hlines}
\toprule
Problem Category & Who to Contact \\
Physical switch hardware failure (broken cable, switch not activating) & AT Specialist or school OT \\
Bluetooth interface malfunction & AT Specialist; manufacturer support if under warranty \\
iPadOS Switch Control behavior not matching documentation & AT Specialist; check iOS version, update may be needed \\
Evermusic blocked by MDM; cannot install VLC either & District IT / MDM Administrator \\
Student behavior / engagement concerns across multiple sessions & BCBA, behavior specialist, or IEP team \\
Visual access concerns (student cannot see screen adequately) & TVI and AT Specialist \\
Motor access concerns (student cannot reliably activate switch) & OT/PT; AT Specialist for switch alternatives \\
Family/home setup not replicating school setup & Share Section~10 of this manual with family; consult AT Specialist for home equipment loan \\
\bottomrule
\end{longtblr}


\textbf{DOCUMENTATION REMINDER}

When a technical problem affects a session, note it in the student's
data sheet: what the problem was, when it occurred, and how it was
resolved.

Recurring technical problems should be flagged at IEP meetings as an
access barrier.

Missing data due to technical issues is \textbf{not} the same as skill
regression, note this clearly in progress reports.


\chapter{Three-switch instruction}\label{section-8-three-switch-extension}
\chaptermark{Three-switch instruction}

\marginnote{CAUTION: Three-switch training requires OT/PT sign-off on the student's physical capacity to access a third switch position. Never introduce Switch~3 without this consultation, adding an unreachable switch creates frustration, not independence.}
\marginnote{Section~8.2 for interface compatibility. Section~6.4 for the three-switch data sheet.}


\textbf{PREREQUISITES, Do NOT proceed without confirming all of the
following:}

\begin{itemize}

\item
  Student has achieved Phase 4 independence with the two-switch system
 , completing 3 full independent scan-select-listen-return cycles per
  session with zero prompts across 3 consecutive sessions
\item
  Student demonstrates understanding of Switch 1 and Switch 2 functions
  when asked to demonstrate or label them
\item
  Student can physically access a third switch position without
  compromising access to existing switches. Consult OT/PT
\item
  IEP team has confirmed that bidirectional scanning is a meaningful
  functional goal for this student
\end{itemize}


\section{Why Add a Previous-Item
Switch?}\label{why-add-a-previous-item-switch}

\marginnote{The same bidirectional logic applies to YouTube playlist navigation (Section~13) and to BARD/Sora book lists (Section~12). Students who master three-switch scanning in music contexts generalize to longer content lists more efficiently.}

\begin{itemize}

\item
  With only a forward switch, if a student moves the highlight past a
  preferred song, they must scan all the way around the full list to
  reach it again
\item
  With a backward switch, the student can immediately back up one step,
  greatly reducing scanning time and frustration on long lists
\item
  For students with fatigue (cerebral palsy, low endurance): reducing
  unnecessary switch presses has direct quality-of-life impact
\item
  For students with visual impairments: being able to back up one step
  to a missed song is meaningful because visually tracking a highlight
  through a long list is difficult
\item
  For students with intellectual disabilities: only introduce the third
  switch if the student has demonstrated clear frustration with
  over-scanning, not as a preemptive upgrade, the added cognitive
  complexity may outweigh the benefit
\end{itemize}


\textbf{THREE SWITCHES ARE NOT INHERENTLY BETTER THAN TWO}

More switches = more complexity. A student who is fully independent and
efficient with two switches is achieving a high-quality outcome. The
third switch is worth adding when it genuinely reduces scanning effort
or frustration for a specific student, not as a standard upgrade or
progression milestone.


\section{Hardware and Interface
Compatibility}\label{hardware-and-interface-compatibility}

Most professional-grade Bluetooth interfaces support at least 3
simultaneous switch inputs. Verify your interface:

\begin{longtblr}{colspec={X[2,l] X[1,l] X[2,l]},rowhead=1,hlines}
\toprule
Interface & Port Count & Three-Switch Compatible? \\
Blue2 Bluetooth Switch & 2 ports & NO, Not compatible without upgrade \\
Tecla Shield & 4 ports & YES, Compatible \\
RJ Cooper Bluetooth Switch Interface & 2-4 ports (varies) & Check your specific model \\
BJOY Stick / Ring & Multiple ports & Verify with documentation \\
AbleNet Cue & 2 ports & NO, Not compatible without upgrade \\
\bottomrule
\end{longtblr}


\textbf{IF YOUR INTERFACE ONLY HAS 2 PORTS}

You will need either a different interface or a switch junction box that
allows multiple switches to be combined into one port. Contact your AT
specialist before purchasing, they can advise on whether a new
interface or adapter is more appropriate and whether district-owned
devices are available for loan.


\textbf{Switch placement principles:}

\begin{longtblr}{colspec={X[1,l] X[1,l]},rowhead=1,hlines}
\toprule
Principle & Rationale \\
Keep Switch~1 (Next) and Switch~2 (Select) in their existing positions & The student has already learned these positions. Moving them resets motor memory \\
Place Switch~3 (Previous) on the opposite side of Switch~1, or in a clearly distinct location & Reduces accidental activation. The student should feel a different spatial position for ``back'' \\
Use a different color or texture for Switch~3 if possible & Supports discrimination, especially for students with visual impairments or intellectual disabilities \\
Test the layout with the student's full arm/hand range before finalizing & Some students with CP can only reliably reach switches in specific arc positions \\
Document the layout with a photo or diagram & Ensures consistency across teachers, paras, substitutes, and home \\
\bottomrule
\end{longtblr}

\section{iPadOS Configuration for Three
Switches}\label{ipados-configuration-for-three-switches}

\begin{enumerate}
\item
  Go to Settings \textgreater{} Accessibility \textgreater{} Switch
  Control \textgreater{} Switches.
\item
  Tap Add New Switch \textgreater{} External (Bluetooth).
\item
  Activate Switch 3 when prompted.
\item
  Name it \textbf{``Back''} or \textbf{``Previous''} and set action to
  \textbf{Move to Previous Item}.
\item
  Tap Save.
\end{enumerate}

\textbf{Verify all three switch assignments:}

\begin{longtblr}{colspec={X[1,l] X[2,l] X[3,l]},rowhead=1,hlines}
\toprule
Switch & Action & How to Verify \\
Switch~1 & Move to Next Item & Press, highlight advances forward by one item \\
Switch~2 & Select Item & Press, highlighted item activates (song plays) \\
Switch~3 & Move to Previous Item & Press, highlight moves backward by one item \\
\bottomrule
\end{longtblr}

All other Switch Control settings remain identical to the two-switch
configuration (Step Scanning, Auto Scanning OFF, Hold Duration 0.5 sec,
Cursor Large, Resume Scanning Automatic).

\section{Three-Switch Training
Sequence}\label{three-switch-training-sequence}

\subsection{Pre-Training: Introduce Switch 3 Concept Before
Sessions}\label{pre-training-introduce-switch-3-concept-before-sessions}

\begin{enumerate}
\item
  While the student is in a non-session moment (calm, not mid-task):
  show them Switch 3. Say: \emph{``This is a new switch. It moves
  backward, it goes back to the last song.''}
\item
  Demonstrate: place the highlight on a song, then press Switch 3. The
  highlight moves back one. Say: \emph{``See, it went back.''}
\item
  Let the student try Switch 3 alone, with no consequence (no song needs
  to play). Just practice the backward movement.
\item
  Do 3-5 practice presses of Switch 3 only, with narration. End. Do not
  combine with a full session yet.
\end{enumerate}

\subsection{Three-Switch Phase 1, Introduction with Physical
Guidance}\label{three-switch-phase-1-introduction-with-physical-guidance}

\textbf{Goal:} Student experiences all three switches in a real session
with full support. Limit to 1-2 Switch 3 uses per session.

\begin{enumerate}
\item
  Run a normal session with Switch 1 and Switch 2 as usual (5 songs,
  familiar content).
\item
  At a natural moment when the student overshoots a preferred song:
  physically guide Switch 3 press. Say: \emph{``Oops, let's go
  back.''} Highlight moves back.
\item
  Then guide Switch 2 to select it. Say: \emph{``There it is, you
  went back and got it.''}
\end{enumerate}

\textbf{Phase 1 criterion:} Student tolerates three-switch session
without increased distress across 2 sessions.

\subsection{Three-Switch Phase 2, Gestural Prompting for Switch
3}\label{three-switch-phase-2-gestural-prompting-for-switch-3}

\begin{enumerate}
\item
  Run normal two-switch scanning. When student overshoots a song:
  \textbf{point to Switch 3}. Wait 3 seconds.
\item
  If student activates Switch 3: reinforce clearly. \emph{``You went
  back! Now pick it.''}
\item
  If no response after 3 seconds: provide light wrist prompt, then back
  off on next trial.
\item
  Do not prompt Switch 3 preemptively. Only prompt when the student has
  clearly missed a preferred song.
\end{enumerate}

\textbf{Phase 2 criterion:} Student activates Switch 3 with gestural
prompt (point only) on \textbf{4 of 5} overshoot opportunities across
\textbf{2 sessions}.

\subsection{Three-Switch Phase 3, Verbal Prompt
Only}\label{three-switch-phase-3-verbal-prompt-only}

\begin{enumerate}
\item
  When student overshoots: say once: \emph{``Did you miss one? You can
  go back.''}
\item
  Do not point. Wait 10 seconds.
\item
  If student uses Switch 3: reinforce. If not: allow them to continue
  scanning forward. That is also acceptable, do not require Switch 3
  use.
\end{enumerate}

\textbf{Phase 3 criterion:} Student spontaneously uses Switch 3 (without
any prompt) at least once per session across \textbf{3 consecutive
sessions} when an overshoot occurs.

\subsection{Three-Switch Phase 4, Independent Three-Switch
Use}\label{three-switch-phase-4-independent-three-switch-use}

No prompts. Student drives all three switches independently. Record data
using Data Sheet C (Section 6.4).

\textbf{Note:} It is acceptable for the student to rarely or never use
Switch 3. If they prefer to scan forward all the way around, that is a
valid strategy. The goal is that Switch 3 is \emph{available} as a tool,
not that it must be used.

\textbf{Phase 4 criterion:} Student demonstrates controlled use of all
three switches across \textbf{5 consecutive sessions} with zero prompts,
regardless of how frequently Switch 3 is used.

\section{Three-Switch
Troubleshooting}\label{three-switch-troubleshooting}

\begin{longtblr}{colspec={X[1,l] X[1,l] X[2,l]},rowhead=1,hlines}
\toprule
Problem & Likely Cause & Fix \\
Student confuses Switch~1 and Switch~3 & Switches look or feel similar; spatial position unclear & Differentiate switches: different color, texture, or size. Adjust spacing. Label with tactile marker \\
Student presses Switch~3 immediately after Switch~2 (goes back after selecting) & Switch~3 in accidental activation position & Review placement with OT. Increase Hold Duration to reduce accidental presses \\
Third switch not recognized by iPadOS & Interface only has 2 ports, or connection not made & Verify interface port count. Check cable seating. Re-pair interface after adding third switch \\
Student refuses to use Switch~3 at all & Preference for existing strategy & Do not force. Continue two-switch sessions. Re-introduce Switch~3 concept in 2-3 weeks \\
Highlight moves in wrong direction with Switch~3 & Switch~3 assigned to Move to Next Item instead of Move to Previous Item & Settings $>$ Accessibility $>$ Switch Control $>$ Switches. Correct the action for Switch~3 \\
Student's scanning efficiency gets WORSE after adding Switch~3 & Cognitive overload from three-choice decision & Return to two-switch system. Wait until two-switch performance is even more solid before re-introducing \\
\bottomrule
\end{longtblr}

\chapter{Disability specific adaptations}\label{section-9-disability-specific-adaptations}
\chaptermark{Disability specific adaptations}

\marginnote{The adaptations in this section apply equally to all app contexts in this manual: Evermusic, VLC, NLS BARD, Bookshare, Sora, and YouTube. Document all adaptations in the IEP as access supports, not informal accommodations.}

Every student is different. The adaptations below are organized by
disability area. Consult with specialists on your team to determine
which apply to your individual student. Adaptations should be documented
in the IEP.

\section{Students with Visual
Impairments}\label{students-with-visual-impairments}

\marginnote{VoiceOver + Switch Control is a powerful combination for blind students: VoiceOver speaks each item as it is highlighted, eliminating the need for teacher narration. Ask your AT specialist to configure this before the first session with a blind student.}
\marginnote{Section~12 for NLS BARD and Bookshare, which are specifically designed for blind and print-disabled users and carry full VoiceOver support.}

\begin{itemize}

\item
  Increase cursor size to Large and use high-contrast color (yellow on
  dark background, or bright red)
\item
  Increase screen brightness to maximum during sessions
\item
  Verbally narrate each song title aloud as the highlight moves over it,
  so the student can track by sound
\item
  Position the iPad at the angle and distance the TVI recommends for the
  student's visual profile
\item
  \textbf{For students who are blind or have very limited usable
  vision:} the switch action itself still teaches cause-and-effect and
  choice-making. Focus on the sound cue (music starting) as the reward
  rather than the visual highlight. The cause-effect chain still works
 , the student hears the song title spoken aloud, presses if they
  want it, hears the song play. The visual highlight is irrelevant for
  these students
\item
  Consider augmenting with a tactile cue: a light tap on the student's
  arm as each song is announced signals that a new option is being
  presented, without physically prompting the switch press
\item
  Consult TVI for whether VoiceOver + Switch Control could automate song
  title narration so the teacher does not need to speak every title
  manually
\item
  For three-switch systems: label all three switches with tactile
  markers (bump dots, raised shapes, different textures) the student can
  identify by touch
\end{itemize}

\section{Students with Cerebral
Palsy}\label{students-with-cerebral-palsy}

\marginnote{CAUTION: Never adjust switch position or mounting between sessions without OT/PT approval. Even a small positional change can make a switch inaccessible to a student with CP and may not become apparent until mid-session.}
\marginnote{Two short sessions per day (5-8 min each) consistently outperform one long session for students with CP-related fatigue. Document session length and fatigue patterns in the IEP as part of the access plan.}

\begin{itemize}

\item
  Use the same switch position and mounting determined by the OT/PT. Do
  not improvise positioning
\item
  If the student activates switches accidentally: increase Hold Duration
  to 0.75-1.0 seconds so they must sustain pressure intentionally
\item
  If pressing is difficult: ask the AT specialist or OT/PT if a
  different switch type (sip-and-puff, head switch, larger/lighter
  switch) would be more appropriate
\item
  Allow extra time. Do not rush. Fatigue is real, shorter sessions
  (5-8 minutes) are better than longer frustrating ones. Two short
  sessions per day may be more productive than one long one
\item
  In the cause-effect stage: involuntary movements may produce
  accidental switch presses. Track timing carefully, presses that
  occur while the highlight is moving (not during a pause) should not
  count as intentional selections
\item
  For three-switch systems: the third switch must not require a movement
  pattern that interferes with Switches 1 and 2. Some students may do
  better with all three switches on a single mounting arm, accessed by
  the same limb using different movements. Others may benefit from foot
  or head switch placement for Switch 3. \textbf{Consult OT/PT before
  finalizing the three-switch layout}
\item
  Celebrate each intentional press, even if it hits the wrong song. The
  movement practice is valuable
\end{itemize}

\section{Students with Intellectual
Disabilities}\label{students-with-intellectual-disabilities}

\marginnote{For students with intellectual disabilities, audiobooks via NLS BARD or Sora (Section~12) may be more motivating than music once two-switch mastery is established, particularly for students who respond well to narrative content or story listening.}

\begin{itemize}

\item
  Keep the song list very short (3-5 songs maximum) and use only songs
  the student already loves
\item
  Use consistent language every session. Pick one phrase and stick with
  it: always \emph{``Find a song,''} not sometimes \emph{``What do you
  want?''}
\item
  Pair the switch press with a clear, immediate reward, the music
  starting. Avoid adding extra steps or praise that might confuse the
  connection
\item
  If the student is not interested in music: work with the teacher to
  find a motivating activity that uses the same two-switch setup (video
  clips, a favorite audiobook, environmental sounds)
\item
  Routine matters. Try to run sessions at the same time and in the same
  place each time
\item
  Extend Phase 0 as needed. Some students may need 10-15 sessions at
  Phase 0 before demonstrating clear discrimination. This is normal. Do
  not rush to Phase 1
\item
  For three-switch systems: only introduce Switch 3 after extended, very
  solid two-switch mastery (consider requiring 5+ consecutive mastery
  sessions rather than 3). Use concrete language: \emph{``Oops switch''}
  or \emph{``Back switch.''} If introducing Switch 3 causes regression
  in Switch 1 and 2 use: discontinue and return to two-switch protocol.
  This is not a failure, it is data
\end{itemize}

\chapter{Home practice guide}\label{section-10-family-and-caregiver-home-practice-guide}
\chaptermark{Home practice guide}

\marginnote{This section is written for families and may be printed as a standalone handout. It covers Evermusic and VLC home use. For families whose child also uses NLS BARD or Sora at school, the teacher should share the relevant app-specific steps from Section~12 as a supplement.}


\textbf{FOR TEACHERS:} This section is written in plain, accessible
language and may be printed separately as a parent handout. Even 5
minutes of home practice a few times a week makes a meaningful
difference in skill acquisition speed.


\section{What Your Child Is
Learning}\label{what-your-child-is-learning}

Your child is working on a skill called switch scanning. They are
learning to use two switches to find and play their favorite music
independently. The music is the motivator, but what they're really
learning is how to make decisions and control what happens around them.

Here is what it looks like step by step:

\begin{longtblr}{colspec={X[1,l] X[7,l]},rowhead=1,hlines}
\toprule
Step & What Happens \\
\textbf{Step 1} & A list of songs appears on the iPad screen \\
\textbf{Step 2} & Child presses Switch~1 (the ``Move'' switch), a colored highlight box moves to the next song \\
\textbf{Step 3} & Child keeps pressing Switch~1 until the highlight is on a song they want \\
\textbf{Step 4} & Child presses Switch~2 (the ``Pick'' switch), the song starts playing \\
\textbf{Step 5} & Song plays all the way through. When it ends, the list comes back on its own \\
\textbf{Step 6} & Child presses Switch~1 again to find and pick the next song, no prompts needed! \\
\bottomrule
\end{longtblr}


\textbf{Remind your child before each session:}

Switch 1, MOVE: \emph{``This switch moves to the next song.''} Switch
2, PICK: \emph{``This switch picks the song and plays it.''}

Tip: If your child's teacher has labeled the switches with color tape or
a symbol, keep those labels on at home too.


\section{Setting Up at Home}\label{setting-up-at-home}

\textbf{What you need:}
\begin{itemize}

\item The iPad your child uses at school (or a home iPad with the same apps)
\item The Bluetooth switch interface and two physical switches (sent home by the teacher)
\item The Evermusic app (or VLC if your district uses that instead)
\item MP3 music files already loaded onto the iPad (your child's teacher can help with this)
\item A quiet space with minimal distractions for the first few sessions
\end{itemize}

\textbf{Before you start, 3 quick checks:}

\begin{longtblr}{colspec={X[1,l] X[2,l] X[2,l]},rowhead=1,hlines}
\toprule
Check & How to Do It & What You Should See \\
Switches are connected & Turn on the Bluetooth switch box. Open iPad Settings $>$ Bluetooth. & A checkmark or ``Connected'' next to the device name \\
Switch scanning is on & Open Settings $>$ Accessibility $>$ Switch Control. & Green toggle at top of Switch Control page \\
Song list is visible & Open Evermusic and tap ``Local Files.'' & A simple numbered or titled list of songs \\
\bottomrule
\end{longtblr}

\textbf{If the switches are not connected:}

\begin{itemize}

\item
  Make sure the Bluetooth switch box is charged and powered on
\item
  On the iPad: Settings \textgreater{} Bluetooth \textgreater{} look for
  the device. If it shows ``Not Connected,'' tap it to reconnect
\item
  If it doesn't appear at all, turn Bluetooth off and back on, then try
  again
\item
  Still not working? Contact your child's teacher
\end{itemize}

\section{Suggested Session Length}\label{suggested-session-length}

\begin{itemize}

\item
  \textbf{Beginners (just starting):} 5-10 minutes, 2-3 song choices
\item
  \textbf{Building confidence:} 10-15 minutes, 4-6 song choices
\item
  \textbf{Independent (rare prompts):} Let your child lead, end when
  they choose to stop
\end{itemize}

\section{How Much Help to Give}\label{how-much-help-to-give}

Your child's teacher will tell you which level of help is right for your
child right now.

\begin{longtblr}{colspec={X[1,l] X[2,l]},rowhead=1,hlines}
\toprule
Help Level & What to Do \\
Full guidance (just starting) & Put your hand gently over their hand and press the switch together. Say the action out loud: \emph{``Move to next.'' / ``Pick that song.''} \\
Light guidance & Rest your hand lightly on their wrist as a reminder, don't push. Point to the switch and wait 3 seconds \\
Pointing only & Point to the switch without touching. Say: \emph{``What switch do you need?''} \\
Words only & Say: \emph{``Find a song.''} Wait. Only speak again if 10 seconds pass with no movement \\
Independent (the goal!) & Say: \emph{``Music time.''} Then stay quiet and let your child drive completely \\
\bottomrule
\end{longtblr}

\section{Between Songs}\label{between-songs}

When a song ends and the list comes back:

\begin{itemize}

\item
  \textbf{Wait silently for at least 10 seconds}
\item
  If your child doesn't start scanning: gently remind them once --
  \emph{``Find another song if you want.''}
\item
  If they look away or seem done: that is okay. Music time is over for
  now
\item
  Do not rush this moment. The pause is part of the learning
\end{itemize}


\textbf{WHAT COUNTS AS SUCCESS}

\begin{itemize}

\item
  Your child pressed a switch independently, even once. That is
  progress.
\item
  Your child tolerated the session without distress. That matters.
\item
  Your child showed any sign of preference, leaning toward a song,
  looking at the screen, a smile when music played.
\end{itemize}

You do NOT need your child to complete the full cycle independently
right away. Partial steps are wins.


\section{Home Troubleshooting}\label{home-troubleshooting}

\begin{longtblr}{colspec={X[1,l] X[1,l] X[1,l]},rowhead=1,hlines}
\toprule
Problem & What It Looks Like & What to Try \\
Switches don't seem to work & Pressing a switch does nothing; no highlight moves & Check Bluetooth (Settings $>$ Bluetooth). Make sure Switch Control is ON. Try turning Bluetooth off and on. \\
Highlight doesn't appear & No colored box around any item & Switch Control may have turned off. Go to Settings $>$ Accessibility $>$ Switch Control and toggle ON. \\
Wrong switch does wrong thing & Move switch picks a song; Pick switch moves highlight & Switches may have been accidentally re-assigned. Contact the teacher to reset. \\
Song plays then another plays automatically & Music keeps playing without child choosing & Open Evermusic $>$ Settings $>$ Player $>$ set Repeat Mode to Stop. Ask teacher if unsure. \\
Child gets frustrated quickly & Distress, pushing away the iPad, refusing to engage & Shorten the session. Reduce song list to 2-3 favorites. Go back to full physical guidance. End on a success, even a small one. \\
Child presses Switch~1 over and over without stopping & Highlight keeps cycling without a selection & This is okay, they are exploring. Wait. After 2-3 cycles, gently say: \emph{``Pick one!''} and prompt Switch~2 if needed. \\
iPad screen goes dark mid-session & Screen turns off while child is working & Increase Auto-Lock: Settings $>$ Display and Brightness $>$ Auto-Lock $>$ set to 5 minutes or Never during sessions. \\
Can't find the song list in Evermusic & App shows cloud storage or asks to sign in & Tap the house/home icon or ``Local Files'' tab. Never sign in to cloud services. Contact teacher if this persists. \\
\bottomrule
\end{longtblr}

\section{Keeping in Touch with the
Teacher}\label{keeping-in-touch-with-the-teacher}

Home practice is most powerful when the teacher knows how it is going.

\begin{longtblr}{colspec={X[1,l] X[1,l]},rowhead=1,hlines}
\toprule
What to Share & Why It Helps \\
How many songs your child chose independently & Helps teacher know if the song list needs to be shorter or longer \\
What level of prompting you used & Lets teacher adjust the phase they are working on at school \\
Any switch or app problems you ran into & Teacher can troubleshoot or send a spare switch box home \\
Songs your child seemed most excited about & Teacher can move preferred songs to the top of the list \\
Whether your child seemed tired or frustrated & May indicate the session was too long, or that positioning needs review \\
\bottomrule
\end{longtblr}


\textbf{YOU ARE PART OF THE TEAM}

Switch access skills take time to build. There is no set timeline and no
``behind.''

What you are doing at home, even in short, imperfect sessions, is
genuinely meaningful for your child's independence and communication.

If you ever feel uncertain about what to do, reach out to the teacher.
There are no wrong questions.


\chapter{Full skill progression quick reference}\label{section-11-full-skill-progression-quick-reference}
\chaptermark{Full skill progression quick reference}

\marginnote{Sections~12 and~13 extend this progression to accessible reading apps and YouTube. Once two-switch independence is established in music contexts, these app contexts become natural IEP generalization targets.}

\section{11.1 Master Progression
Table}\label{master-progression-table}

\begin{longtblr}{colspec={X[2,l] X[1,l] X[3,l] X[1,l]},rowhead=1,hlines}
\toprule
Stage & Section & Mastery Criterion & Move Forward When\ldots{} \\
Cause-Effect Phase 0 & Section~3 & 3/5 presses in pause window across 2 sessions, no physical prompt & Student responds in pause window consistently without physical prompt \\
Cause-Effect Phase 1 & Section~3 & 4/5 preferred pressed; $\leq$1 non-preferred pressed per session $\times$ 3 sessions & Clear discrimination between preferred and non-preferred demonstrated \\
Cause-Effect Phase 2 & Section~3 & 4/5 preferred with 3-sec pause, no verbal narration $\times$ 3 sessions & Student selects preferred items without verbal support \\
Two-Switch Phase 1 & Section~4 & Tolerates physical guidance $\times$ 2 consecutive sessions & Student accepts guidance without resistance \\
Two-Switch Phase 2 & Section~4 & Activates each switch with gestural prompt on 4/5 trials & Consistent switch use with point prompt only \\
Two-Switch Phase 3 & Section~4 & Full scan-select-listen-return cycle with verbal only on 4/5 trials $\times$ 2 sessions & Re-initiation after song end with verbal cue only \\
Two-Switch Phase 4 & Section~4 & 3 full independent cycles with zero prompts $\times$ 3 consecutive sessions & All three components of mastery met \\
Generalization & Section~4.8 & Mastery across 2+ different settings or partners & Skill is truly functional and context-independent \\
Three-Switch Introduction & Section~8 & Phase~4 two-switch mastery confirmed; OT/PT and IEP team consulted & Bidirectional scanning has clear functional benefit for the student \\
Three-Switch Phase 4 & Section~8 & Controlled three-switch use $\times$ 5 consecutive sessions, zero prompts & Scanning efficiency is maintained or improved with three switches \\
Reading Apps (BARD / Bookshare / Sora) & Section~12 & Phase~2 two-switch independence in music context & Student ready to generalize scan-select skill to reading app context \\
YouTube Navigation & Section~13 & Phase~3 or Phase~4 two-switch independence & Student ready to apply scanning to longer-form video content \\
\bottomrule
\end{longtblr}

\section{11.2 IEP Integration Notes}\label{iep-integration-notes}

\begin{itemize}

\item
  Data from all sessions should be available at every IEP meeting.
  Present graphs or tables when possible
\item
  Recurring technical problems should be documented as \textbf{access
  barriers} in the IEP
\item
  Phase level and prompt hierarchy should be documented in the IEP as
  the current instructional level
\item
  Generalization targets should be listed as IEP goals once Phase 4
  mastery is established
\item
  Missing data due to technical issues is \textbf{not} the same as skill
  regression, note this clearly in progress reports
\end{itemize}

\section{11.3 When Progress Stalls}\label{when-progress-stalls}

If a student has not met the mastery criterion for a phase after 10+
sessions, do not simply continue. Investigate:

\begin{itemize}

\item
  Are the songs actually preferred? Get fresh preference data from the
  family
\item
  Can the student physically access the switch reliably? Consult OT/PT
\item
  Is the pause window (in cause-effect stage) long enough? Try extending
  to 7-10 seconds
\item
  Is the cause-effect concept established in other contexts? Try a
  simpler single-switch toy
\item
  Is there a sensory barrier? Can the student see the screen and hear
  the music adequately?
\item
  Is the session environment stable and low-distraction?
\item
  Is the \textbf{same session structure} being used consistently across
  all staff? Staff-to-staff variation in prompting can stall progress
\end{itemize}

\section{11.4 How This Protocol Fits the Full
Progression}\label{how-this-protocol-fits-the-full-progression}

\marginnote{Sections~12 and~13 extend the protocol into accessible reading and video contexts.}

\begin{longtblr}{colspec={X[1,l] X[2,l]},rowhead=1,hlines}
\toprule
Document / Section & Stage \\
\textbf{Section 3 (this manual)} & One switch: cause-and-effect; student selects only \\
\textbf{Sections 4-6 (this manual)} & Two switches: student navigates and selects independently (Evermusic) \\
\textbf{Section 5 (this manual)} & Two switches: VLC variant for MDM-blocked environments \\
\textbf{Section 8 (this manual)} & Three switches: adding Move to Previous Item for students ready for bidirectional scanning \\
\textbf{Section 12 (this manual)} & Two or three switches: NLS BARD, Bookshare, and Sora accessible reading app setup and navigation for students with print disabilities and/or visual impairments \\
\textbf{Section 13 (this manual)} & Two or three switches: YouTube navigation for video-based learning, generalization of scanning skills to an open web environment \\
\bottomrule
\end{longtblr}

Each stage has explicit mastery criteria. Students should not move
forward until criteria are met.


\textbf{FINAL NOTE}

This protocol is a framework, not a script. Every student is different,
and the best teachers adapt thoughtfully. The principles, use the
least intrusive prompt, collect data on every session, let the data
drive decisions, and prioritize genuine motivation, apply regardless
of what adaptations you make.

When in doubt: do less. Wait longer. Let the student show you what they
can do.


\emph{Switch Access Training Protocol Series, Compiled from: Switch
Access Protocol v2, Appendix A (Parent Guide), Appendix B
(Troubleshooting), Appendix C (Three-Switch System), Appendix D (Cause
and Effect). Please share feedback or questions with your AT specialist
or protocol coordinator.}

\chapter{Accessible reading apps}\label{section-12-accessible-reading-apps}
\chaptermark{Accessible reading apps}

\textbf{PREREQUISITE FOR THIS SECTION}

Students should have achieved at minimum Phase~2 two-switch independence
(gestural-prompt level) in music contexts (Sections~4-5) before
beginning accessible reading app instruction. The scan-select cognitive
framework must be established before extending it to a new app context.

\marginnote{This section mirrors the structure of Sections~2, 4, and~5. The same switch hardware, Switch Control configuration, and prompting hierarchy apply. What changes is the app, the content type, and the navigation structure.}

\section{12.1 Overview and App
Comparison}\label{accessible-reading-app-overview}

\marginnote{CAUTION: NLS BARD requires registration through the National Library Service. Students must qualify as print-disabled (including visual impairments and physical disabilities that prevent holding a book or turning pages). Contact your school's TVI or AT specialist to begin the application process.}

NLS BARD, Bookshare, and Sora serve overlapping but distinct purposes
for students who are blind, have low vision, or have physical
disabilities that make print inaccessible. All three support switch
access via iPadOS Switch Control and work with the same Bluetooth
interface used for Evermusic and VLC.

\begin{longtblr}{colspec={X[1,l] X[2,l] X[3,l] X[2,l] X[1,l] X[2,l]},rowhead=1,hlines}
\toprule
App & Publisher / Source & Content Type & Eligibility & Cost & MDM Approval \\
NLS BARD Mobile & National Library Service (Library of Congress) & Audiobooks, braille files, magazines, federally funded & Print-disabled users (VI, physical disability) & Free & Usually approved, federal app \\
Bookshare & Benetech (nonprofit) & Ebooks, audiobooks, DAISY, braille-ready & Students with print disabilities; IEP documentation required & Free for qualifying students & Usually approved, nonprofit educational \\
Sora & OverDrive / Libby (school library) & School/library ebooks and audiobooks & Any student at a participating school & Free through school library & Varies, confirm with IT \\
\bottomrule
\end{longtblr}

\marginnote{Section~2.1 for equipment requirements. The same Bluetooth interface supports all apps in this section.}

\section{12.2 Why These Apps Matter for Switch
Users}\label{why-accessible-reading-apps-matter}

\marginnote{For students who are blind or have low vision, NLS BARD and Bookshare are not simply ``nice to have'', they provide access to the same curriculum and recreational reading available to sighted peers. Configuring switch access to these apps is an equity issue, not an optional enhancement.}

\begin{itemize}

\item
  All three apps support \textbf{VoiceOver + Switch Control}: as the
  highlight scans each on-screen item, VoiceOver speaks its name or
  content aloud, eliminating the need for teacher narration and
  giving blind students fully independent access to book and title
  selection
\item
  All three offer \textbf{audio playback} (audiobooks, read-aloud), so
  students who cannot access visual text receive full content access
  through sound
\item
  Books and audio are downloaded \textbf{directly within each app}. No
  USB transfer is required, a key advantage over the MP3 workflow in
  Sections~2 and~5
\item
  For students with \textbf{combined visual and physical impairments},
  these apps provide a meaningful reading context that music alone
  cannot: narrative content, curriculum-aligned titles, and library
  books their peers are reading
\end{itemize}

\section{12.3 NLS BARD Mobile, Setup and Switch
Configuration}\label{nls-bard-setup}

\marginnote{NLS BARD registration must be completed by a qualified certifier (TVI, AT specialist, or librarian). Allow 1-4 weeks for account approval. Begin the process well before you plan to start using the app.}

\textbf{Step 1, Register for NLS BARD}

\begin{enumerate}
\item
  Contact the student's TVI or AT specialist to complete the NLS BARD
  eligibility application. Qualifying conditions include visual
  impairments, and physical disabilities that prevent holding a book or
  turning pages (including CP and other orthopedic impairments).
\item
  Once approved, the student receives an NLS BARD account (username and
  password). Record these in the student's AT documentation.
\item
  Download \textbf{NLS BARD Mobile} from the App Store. The developer
  is the Library of Congress / NLS.
\end{enumerate}

\textbf{Step 2, Initial App Configuration (Teacher, One-Time)}

\begin{enumerate}
\item
  Open BARD Mobile and log in with the student's NLS credentials.
\item
  Download 3-5 books or magazines the student is known to prefer.
  Navigate: \textbf{Bookshelf \textgreater{} Browse/Search} to find
  titles; tap a title and select \textbf{Download}.
\item
  Once downloaded, titles appear in \textbf{My Books} (the bookshelf
  view). This is the list the student will scan.
\item
  \textbf{Set playback speed:} within any open audiobook, tap the
  speed icon and set to 1.0x (normal) as a starting point. Adjust
  after observing the student's comprehension.
\item
  \textbf{Turn off auto-advance between chapters:} In BARD settings,
  confirm that playback pauses at chapter end rather than advancing
  automatically, this mirrors the Evermusic Repeat Stop philosophy
  and gives the student deliberate control over when to continue.
\item
  Leave the app on the \textbf{My Books} shelf (the book list view)
  before each session.
\end{enumerate}

\marginnote{Limit the bookshelf to 5-8 titles for early sessions. A shorter scan list reduces the number of Switch~1 presses to reach a desired book, directly lowering motor demand for students with CP.}

\textbf{Step 3, Switch Control Verification for BARD}

\begin{enumerate}
\item
  Enable Switch Control (Settings \textgreater{} Accessibility
  \textgreater{} Switch Control, green toggle). Confirm Switch~1
  advances the highlight through the book list and Switch~2 activates
  the highlighted title.
\item
  \textbf{If using VoiceOver + Switch Control:} enable VoiceOver
  (Settings \textgreater{} Accessibility \textgreater{} VoiceOver --
  green toggle). With both active, pressing Switch~1 will move the
  highlight \emph{and} cause VoiceOver to speak the book title aloud.
  This is the recommended configuration for blind or low-vision
  students.
\item
  Test the full cycle before the student's first session: scan to a
  book title (Switch~1), select it (Switch~2), confirm the audiobook
  begins playing.
\end{enumerate}

\marginnote{\textbf{VoiceOver + Switch Control} is the gold standard for blind students. With this combination, every book title, button, and navigation element is spoken as the highlight lands on it, fully autonomous access with no teacher narration needed. Consult the TVI to configure VoiceOver settings appropriately for the student's needs.}

\section{12.4 BARD Session Workflow (6-Step Cycle)}\label{bard-session-workflow}

\begin{longtblr}{colspec={X[1,l] X[7,l]},rowhead=1,hlines}
\toprule
Step & Action \\
\textbf{Step 1} & BARD is open on the My Books shelf. Book list is visible (or spoken via VoiceOver) \\
\textbf{Step 2} & Student presses Switch~1 (Next Item), highlight moves to next book title (VoiceOver speaks the title if enabled) \\
\textbf{Step 3} & Student continues pressing Switch~1 until the desired title is highlighted \\
\textbf{Step 4} & Student presses Switch~2 (Select), audiobook opens and begins playing \\
\textbf{Step 5} & Student listens. At a natural pause point (end of chapter, or teacher-defined time), student presses Switch~2 to pause playback \\
\textbf{Step 6} & Student scans back to the bookshelf button (top-left of player screen) and selects it to return to the list, or continues listening \\
\bottomrule
\end{longtblr}

\marginnote{CAUTION: The BARD player screen has several controls that Switch Control will scan: playback speed, rewind, bookmark, and navigation buttons. In early sessions, provide verbal guidance for which control to select. As with VLC, consider teaching the ``return to bookshelf'' step as a discrete skill before full independent sessions (see Section~5.5 for the model procedure).}

\section{12.5 BARD Pre-Session
Checklist}\label{bard-pre-session-checklist}

\begin{itemize}

\item[\checkbox] 
Bluetooth switch interface is powered on and connected to iPad
\item[\checkbox] 
Switch Control is enabled (Settings \textgreater{} Accessibility)
\item[\checkbox] 
Switch~1 = Move to Next Item confirmed (test: highlight advances)
\item[\checkbox] 
Switch~2 = Select Item confirmed (test: highlighted item activates)
\item[\checkbox] 
VoiceOver is enabled if student has a visual impairment (confirm with
  TVI)
\item[\checkbox] 
BARD is open and showing the My Books shelf (not a book's interior
  or a login screen)
\item[\checkbox] 
Book list contains 5-8 preferred/appropriate titles
\item[\checkbox] 
Playback speed is set appropriately for this student
\item[\checkbox] 
iPad is charged; Auto-Lock set to 5 minutes or Never
\item[\checkbox] 
Data collection sheet is ready
\end{itemize}

\section{12.6 Bookshare App, Setup and Switch
Configuration}\label{bookshare-setup}

\marginnote{Bookshare accounts for qualifying students are free. To set up a student account, visit \texttt{bookshare.org} and follow the Individual Membership process. IEP documentation of the qualifying disability is required. Your AT specialist or TVI can assist.}

\textbf{Step 1, Create a Bookshare Account}

\begin{enumerate}
\item
  Visit \texttt{bookshare.org} and apply for a student Individual
  Membership. The qualifying disability must be documented (visual
  impairment, physical disability, or other print disability). Your
  school's TVI or AT specialist can certify eligibility.
\item
  Once approved (typically within a few days), the student receives an
  email login. Record credentials in AT documentation.
\item
  Download \textbf{Bookshare} (also called the Bookshare Reading App)
  from the App Store.
\end{enumerate}

\textbf{Step 2, App Configuration (Teacher, One-Time)}

\begin{enumerate}
\item
  Log in with the student's Bookshare credentials.
\item
  Search for and download 3-5 titles of known interest to the student.
  Navigate: \textbf{Search or Browse \textgreater{} tap a title
  \textgreater{} Download}. Downloaded books appear in
  \textbf{My Books}.
\item
  Set reading format: for students with visual impairments or who are
  blind, select \textbf{Audio} (TTS or human-narrated) rather than
  visual text. Tap the book cover \textgreater{} settings \textgreater{}
  select audio format.
\item
  Set font size and contrast (for low-vision students): open a book
  \textgreater{} tap the display icon \textgreater{} increase font size
  and select high-contrast mode.
\item
  Leave the app on the \textbf{My Books} shelf before each session.
\end{enumerate}

\marginnote{Bookshare's audio reading format uses text-to-speech. For students who strongly prefer human narration, NLS BARD offers more human-narrated audiobooks. Many Bookshare titles are also available in DAISY format, which supports navigation by chapter, section, and page, more sophisticated navigation that can become a generalization goal once basic switch scanning is mastered.}

\textbf{Step 3, Switch Control Verification for Bookshare}

Same procedure as BARD (Section~12.3, Step~3). The Switch Control
settings are identical. Verify Switch~1 advances highlight through the
book list and Switch~2 opens a title.

\marginnote{Section~12.3 Step~3 for the full Switch Control verification procedure.}

\section{12.7 Sora App, Setup and Switch
Configuration}\label{sora-setup}

\marginnote{Sora connects directly to a school or public library's digital collection. Any student at a participating school can use Sora without a disability certification, making it the most immediately accessible of the three apps for general switch-using students.}

\textbf{Step 1, Verify School Participation and Log In}

\begin{enumerate}
\item
  Confirm your school or district uses OverDrive/Sora for its digital
  library. Ask your school librarian or IT department.
\item
  Download \textbf{Sora} from the App Store (developer: OverDrive).
\item
  Open Sora. Tap \textbf{Find My School} and search for your school or
  district by name or zip code.
\item
  Log in using the student's school credentials (usually the same
  username/password used for school Google or Microsoft accounts).
\end{enumerate}

\textbf{Step 2, App Configuration (Teacher, One-Time)}

\begin{enumerate}
\item
  Browse and borrow 3-5 titles relevant to the student's interests or
  curriculum. Navigate: \textbf{Explore \textgreater{} search or browse
  \textgreater{} tap a title \textgreater{} Borrow}. Borrowed titles
  appear on the student's \textbf{Shelf}.
\item
  For students with visual impairments: select audiobook format when
  borrowing. Tap a title \textgreater{} if both ebook and audiobook
  formats are available, choose \textbf{Borrow Audiobook}.
\item
  Set text display for low-vision students: open an ebook in Sora
  \textgreater{} tap the font/display icon \textgreater{} increase font
  size, adjust line spacing, select a high-contrast theme.
\item
  Leave the app on the \textbf{Shelf} (borrowed books list) before each
  session.
\end{enumerate}

\marginnote{CAUTION: Sora borrows have due dates (typically 7-21 days). Build a routine of checking borrowed titles weekly and renewing or replacing them before they expire. An expired book disappears from the Shelf mid-session, which can be disorienting for the student.}

\textbf{Step 3, Switch Control Verification for Sora}

\begin{enumerate}
\item
  With Switch Control enabled, open Sora to the Shelf. Press Switch~1
  to confirm the highlight advances through borrowed titles.
\item
  Press Switch~2 on a highlighted title to confirm the book or
  audiobook opens.
\item
  For audiobooks: confirm Switch~2 activates the play button once the
  book is open.
\item
  \textbf{Note:} Sora's bottom navigation bar (Shelf, Explore,
  Notifications, etc.) will be included in the Switch Control scan path.
  In early sessions, this may confuse students who accidentally scan
  into the navigation bar. Provide verbal guidance: \emph{``Keep going
 , find your book.''} Over time, teach the student to recognize and
  skip past navigation bar items.
\end{enumerate}

\marginnote{Section~7.2 for general Switch Control scanning problems that may apply to any app, including Sora.}

\section{12.8 Sora Session Workflow (6-Step
Cycle)}\label{sora-session-workflow}

\marginnote{The Sora workflow is nearly identical to the BARD workflow (Section~12.4). Once a student is proficient in one, transition to the other with minimal re-teaching.}

\begin{longtblr}{colspec={X[1.4,l] X[9.8,l]},rowhead=1,hlines}
\toprule
Step & Action \\
\textbf{Step 1} & Sora is open on the Shelf (borrowed titles list) \\
\textbf{Step 2} & Student presses Switch~1, highlight moves to next title (VoiceOver speaks title if enabled) \\
\textbf{Step 3} & Student presses Switch~1 until desired title is highlighted \\
\textbf{Step 4} & Student presses Switch~2, book or audiobook opens \\
\textbf{Step 5} & For audiobooks: student presses Switch~2 on the Play button to begin listening \\
\textbf{Step 6} & Student listens. At natural pause, student scans to the back button (top-left) and returns to the Shelf if desired \\
\bottomrule
\end{longtblr}

\section{12.9 Troubleshooting, Accessible Reading
Apps}\label{accessible-reading-app-troubleshooting}

\marginnote{Section~7 for general Switch Control and hardware troubleshooting that applies to all apps.}

\begin{longtblr}{colspec={X[1,l] X[1,l] X[2,l]},rowhead=1,hlines}
\toprule
Problem & Likely Cause & Fix \\
App shows login screen at start of session & Session timed out or app was force-closed & Log back in with saved credentials. Enable ``Remember me'' if available. For BARD and Bookshare: contact AT specialist if login is a repeated problem, MDM may be logging students out \\
Highlight scans into navigation bar instead of book list & App navigation bar is part of the Switch Control scan path & In early sessions, provide verbal guidance. Over time, teach student to recognize navigation bar items and continue past them. Consider Switch Control's Group Items feature, consult AT specialist \\
VoiceOver speaking too fast for student to follow & Default VoiceOver speech rate too high & Settings $>$ Accessibility $>$ VoiceOver $>$ Speech Rate. Reduce to comfortable pace. Student should hear the title spoken before the highlight moves on \\
Audiobook does not start playing after Switch~2 selection & Highlight may have selected the book cover rather than the Play button & Open the book manually to identify the Play button location. Practice isolated Switch~2 presses on the Play button before full sessions \\
Sora book expired and disappeared from shelf & Loan period ended & Renew or re-borrow the title. Check due dates weekly and build this into teacher or AT technician routine \\
BARD shows no downloaded books (empty shelf) & Books not yet downloaded, or storage cleared & Download books while on Wi-Fi before the session. Store 3-5 books offline so sessions can proceed without internet \\
Bookshare audio format not available for a title & Not all Bookshare titles have audio & Filter search by ``Audio'' format, or use DAISY format with a compatible reader. Contact TVI for format guidance \\
\bottomrule
\end{longtblr}

\section{12.10 Disability-Specific Considerations for Reading
Apps}\label{reading-app-disability-considerations}

\marginnote{For students with combined CP and visual impairments, the combination of VoiceOver + Switch Control + audiobook content is typically the most accessible approach: no visual tracking required, no reading required, and switch access is fully autonomous.}

\textbf{Students with visual impairments / blindness:}
\begin{itemize}

\item
  Enable VoiceOver + Switch Control for full autonomous access to all
  three apps. Verify with TVI that VoiceOver is configured for the
  student's reading level (speech rate, language, punctuation settings)
\item
  Choose audiobook or audio TTS formats rather than ebook visual text
\item
  NLS BARD is the most appropriate starting point for blind students
  due to its human-narrated collection, federal accessibility mandate,
  and zero-cost status
\item
  Tactile labels on Switch~1 and Switch~2 remain important, the
  student cannot see the highlight and relies entirely on switch
  position and VoiceOver feedback
\end{itemize}

\textbf{Students with cerebral palsy / orthopedic impairments:}
\begin{itemize}

\item
  These apps require the same two-switch navigation as Evermusic. No
  additional motor demands are introduced
\item
  Session fatigue consideration: audiobooks allow longer passive
  listening periods between switch interactions, which may reduce
  motor fatigue compared to music (which requires re-selection every
  3-5 minutes)
\item
  For students who struggle with the return-to-bookshelf step: allow
  the teacher to return the app to the shelf between songs and focus
  instruction entirely on the scan-select cycle
\end{itemize}

\textbf{Students with both VI and CP:}
\begin{itemize}

\item
  VoiceOver + Switch Control + audiobook is the recommended
  combination. This eliminates all visual and manual demands beyond
  switch pressing
\item
  Consult both TVI and OT/PT before the first session. Optimal switch
  position and VoiceOver settings should be determined collaboratively
\item
  Document the full configuration (switch positions, VoiceOver
  settings, app, content format) in the IEP as an access support
\end{itemize}

% -------------------------------------------------------------

\chapter{YouTube navigation}\label{section-13-youtube-navigation}

\chaptermark{YouTube navigation}

\marginnote{YouTube is a naturally motivating context for many students and is an excellent generalization target after Evermusic/VLC mastery. The same two-switch workflow applies; the added complexity is managing YouTube's larger, more variable interface.}


\textbf{PREREQUISITE FOR THIS SECTION}

Students should have achieved \textbf{Phase~3 or Phase~4 two-switch
independence} (verbal-prompt or independent level) in music or reading
app contexts before beginning YouTube switch navigation. YouTube's
interface is more complex than Evermusic, VLC, or BARD, and the scan
path is longer. Students not yet at Phase~3 will experience high error
rates that undermine motivation.


\section{13.1 Why YouTube and Important Considerations}\label{why-youtube}

\marginnote{CAUTION: YouTube contains unpredictable content. Before any student uses YouTube via switch access, the teacher must configure a curated, restricted playlist or use YouTube Kids. Standard YouTube should not be used without supervision or content controls. See Section~13.3 for setup options.}

\textbf{Why YouTube is worth including in switch access training:}

\begin{itemize}

\item
  YouTube is one of the most widely used recreational and educational
  media platforms, access to it represents genuine social and
  cultural inclusion
\item
  For students with visual impairments, YouTube provides access to
  described video content (Audio Descriptions), educational channels,
  and accessible media
\item
  YouTube navigation via switch scanning is a meaningful generalization
  of the scan-select-play skill chain established with music apps
\item
  For older students, independently choosing and playing a YouTube
  video is a functional life skill with real-world relevance beyond
  school
\item
  Many curriculum-aligned educational videos (science demonstrations,
  read-alouds, music performances, adaptive PE instruction) are
  available on YouTube, making it a legitimate instructional tool
\end{itemize}

\textbf{Important limitations to address before setup:}

\begin{itemize}

\item
  YouTube's interface has many more scan items than Evermusic or BARD, thumbnails, ads, recommended videos, navigation menus, and
  in-video controls all appear in the scan path
\item
  YouTube's interface changes with app updates; scan paths described
  here may vary slightly between app versions
\item
  Advertisements appear during playback and introduce unexpected scan
  items. Use YouTube Premium (ad-free) if available, or configure
  Restricted Mode to reduce ad frequency
\item
  Internet connectivity is required for standard YouTube. Ensure the
  iPad has reliable Wi-Fi during sessions
\end{itemize}

\section{13.2 Equipment and MDM
Considerations}\label{youtube-equipment-mdm}

\marginnote{TIP: Many districts block YouTube via MDM. If YouTube is blocked on the student's device, contact IT to request an exception for the specific student on accessibility grounds. Document the request in the IEP as an access need. YouTube Kids is more commonly approved and may be an appropriate starting point.}

\begin{itemize}

\item
  All switch hardware, Bluetooth interface, and Switch Control
  configuration is identical to Sections~2 and~4, no additional
  equipment is needed
\item
  The YouTube app (free, Apple App Store) is preferred over the
  browser-based version for switch access, the app interface has
  a more predictable scan path than Safari/Chrome
\item
  \textbf{YouTube Kids} (separate app) is an alternative for younger
  students or when content restriction is required. It has a simpler
  interface and age-appropriate content controls
\item
  Confirm with your district IT/MDM administrator whether YouTube or
  YouTube Kids is permitted on student devices before planning
  instruction
\end{itemize}

\section{13.3 Content Setup, Playlists and Restricted
Mode}\label{youtube-content-setup}

\marginnote{TIP: Creating a teacher-curated playlist is the single most important preparation step for YouTube switch access. A short playlist (5-10 videos) creates a familiar, predictable scan list analogous to the Evermusic song list. Students know what to expect, scan efficiently, and encounter no unexpected content.}

\textbf{Option A, Teacher-Curated Playlist (Recommended)}

\begin{enumerate}
\item
  Using a teacher or school YouTube account (not the student's account),
  create a playlist of 5-10 videos appropriate for the student.
  Suggested starting content: music videos of the student's preferred
  artists, educational content aligned to classroom topics, or
  read-aloud channels.
\item
  Name the playlist clearly (e.g., ``[Student Name] Favorites'').
  Set it to \textbf{Unlisted} or \textbf{Private} for privacy.
\item
  On the student's iPad: log into the YouTube app with the teacher
  account, or share the playlist link and open it in the app.
\item
  Navigate to the playlist and leave the app on the \textbf{playlist
  view} before each session. The playlist view shows a vertical list of
  video thumbnails, the closest analog to the Evermusic song list.
\end{enumerate}

\textbf{Option B, YouTube Kids App}

\begin{enumerate}
\item
  Download the \textbf{YouTube Kids} app from the App Store.
\item
  Set up a child profile with an appropriate content level
  (Preschool, Younger, Older).
\item
  Use the parent controls to select approved channels and block
  inappropriate content categories.
\item
  Leave the app on the \textbf{home screen} or a specific channel
  view before each session.
\end{enumerate}

\textbf{Enabling Restricted Mode (YouTube app):}

\begin{enumerate}
\item
  Open the YouTube app \textgreater{} tap the profile icon (top right)
  \textgreater{} Settings \textgreater{} General \textgreater{}
  Restricted Mode. Toggle ON.
\item
  Note: Restricted Mode is not a substitute for supervision or
  curated playlists, but it provides an additional content filter.
\end{enumerate}

\marginnote{Restricted Mode reduces but does not eliminate all inappropriate content. For students using YouTube independently, a curated playlist (Option A) or YouTube Kids (Option B) provides stronger content control than Restricted Mode alone.}

\section{13.4 iPadOS Switch Control Settings for
YouTube}\label{youtube-switch-control-settings}

\marginnote{Section~2.3 for the complete Switch Control configuration. Settings are identical for YouTube; this subsection covers YouTube-specific considerations only.}

Switch Control settings remain identical to the standard configuration
(Step Scanning, Auto Scanning OFF, Switch~1 = Next Item, Switch~2 =
Select, Cursor Large, Ignore Repeat ON). The following YouTube-specific
adjustments are recommended:

\begin{longtblr}{colspec={X[1,l] X[1,l] X[2,l]},rowhead=1,hlines}
\toprule
Setting & Recommendation & Reason \\
Group Items & Consider enabling for YouTube & Groups sets of related items (e.g., a video thumbnail and its title) so they scan as one unit, reducing total presses \\
Hold Duration & 0.75-1.0 sec for YouTube & YouTube has more items to scan; longer Hold Duration reduces accidental activations on unintended items \\
Cursor Color & High-visibility (yellow or red) & YouTube dark interface makes a bright cursor easier to track \\
\bottomrule
\end{longtblr}

\marginnote{If the student is struggling with the length of the YouTube scan path, ask your AT specialist about configuring Switch Control's \textbf{Group Items} feature. This can significantly reduce the number of presses needed to reach a specific video on the playlist.}

\section{13.5 YouTube Session Workflow, Playlist
Navigation}\label{youtube-session-workflow}

\marginnote{This 7-step workflow is intentionally parallel to the Evermusic workflow (Section~4.1). If you have taught the music protocol, use that language when introducing YouTube: ``It works just like your music app, scan to find a video, pick it, watch it, then scan again.''}

\begin{longtblr}{colspec={X[1.6,l] X[10.6,l]},rowhead=1,hlines}
\toprule
Step & Action \\
\textbf{Step 1} & YouTube app is open on the curated playlist view. Video thumbnails are visible as a vertical list \\
\textbf{Step 2} & Student presses Switch~1 (Next Item), highlight moves to the next video thumbnail \\
\textbf{Step 3} & Student continues pressing Switch~1 until the desired video is highlighted \\
\textbf{Step 4} & Student presses Switch~2 (Select), video opens and begins playing \\
\textbf{Step 5} & Video plays. Student watches or listens. Teacher leaves student without prompting \\
\textbf{Step 6} & When video ends (or student wants to change): student presses Switch~1 to scan to the back/playlist button (top-left of video player) and presses Switch~2 to return to the playlist \\
\textbf{Step 7} & Playlist view reappears. Student repeats from Step~2 \\
\bottomrule
\end{longtblr}

\marginnote{CAUTION: YouTube's autoplay feature will queue the next video automatically if not disabled. This bypasses the student's scanning and selection, exactly the problem that Evermusic's Repeat Stop setting prevents. Disable autoplay before every session: from the video player, tap the autoplay toggle (usually visible at the bottom of the screen) to OFF.}

\section{13.6 Teaching the Return-to-Playlist Step (YouTube
Only)}\label{youtube-return-to-playlist}

Like VLC (Section~5.5), YouTube requires the student to navigate back
to the playlist after a video ends. This step should be taught as a
discrete skill before full independent sessions begin.

\marginnote{Section~5.5 for the VLC return-to-list procedure, which serves as the model for this step.}

\textbf{Isolated Practice Procedure:}

\begin{enumerate}
\item
  Teacher plays a short video clip manually (selects it without switch
  input). Pauses it or lets it finish.
\item
  Present the video player screen to the student.
\item
  Prompt: \emph{``The video is done. Find the list.''} Use the same
  prompting hierarchy from Section~4.2.
\item
  Target: student independently scans to the back/playlist button
  (top-left arrow or playlist icon) and selects it within 20 seconds.
\item
  \textbf{Criterion before full independent trials:} 4 of 5 correct
  returns with verbal prompt only, across 2 sessions.
\end{enumerate}

\section{13.7 Disability-Specific Adaptations for
YouTube}\label{youtube-disability-adaptations}

\textbf{Students with visual impairments / low vision:}
\begin{itemize}

\item
  Enable VoiceOver + Switch Control. With both active, VoiceOver will
  speak each video's title as the highlight lands on it, the student
  hears the title rather than seeing the thumbnail
\item
  Choose videos with \textbf{Audio Description (AD)} tracks when
  available. Channels specifically producing accessible content (e.g.,
  described documentary or educational channels) are preferred
\item
  Set YouTube's text size to Large via Display and Text Size in
  Accessibility settings
\item
  For blind students: the audio of most YouTube content is fully
  accessible. The goal is autonomous \emph{selection} of what to
  watch, VoiceOver + Switch Control provides this
\end{itemize}

\marginnote{YouTube's built-in closed captions and audio descriptions are accessible via Switch Control. As a generalization goal, teach the student to navigate to the CC/AD button during playback. This is a higher-level skill appropriate for students at Phase~4 independence who are ready for more complex in-video navigation.}

\textbf{Students with cerebral palsy / orthopedic impairments:}
\begin{itemize}

\item
  YouTube's longer scan path means more Switch~1 presses to reach a
  desired video. Use the curated playlist (Section~13.3, Option A)
  to keep the list short and reduce motor demand
\item
  Increase Hold Duration to 0.75-1.0 seconds to reduce accidental
  video selections while scanning
\item
  Consider enabling the three-switch system (Section~8) once
  two-switch YouTube use is established, so the student can back up
  if they overshoot a preferred video
\item
  Fatigue monitoring: YouTube sessions may be longer (videos are
  typically longer than songs), but the number of switch presses per
  session is lower. Monitor and document fatigue patterns
\end{itemize}

\section{13.8 YouTube Pre-Session
Checklist}\label{youtube-pre-session-checklist}

\begin{itemize}

\item[\checkbox] 
Bluetooth switch interface is powered on and connected to iPad
\item[\checkbox] 
Switch Control is enabled (Settings \textgreater{} Accessibility)
\item[\checkbox] 
Switch~1 = Move to Next Item confirmed (test: highlight advances)
\item[\checkbox] 
Switch~2 = Select Item confirmed (test: highlighted item activates)
\item[\checkbox] 
VoiceOver is enabled if student has a visual impairment (confirm
  with TVI)
\item[\checkbox] 
YouTube app is open on the curated playlist view (not the home
  feed or search screen)
\item[\checkbox] 
Autoplay is \textbf{OFF} (check in video player before handing
  iPad to student)
\item[\checkbox] 
Restricted Mode is ON, or YouTube Kids is being used
\item[\checkbox] 
Playlist contains 5-10 appropriate videos for this student
\item[\checkbox] 
iPad has Wi-Fi connectivity confirmed
\item[\checkbox] 
Auto-Lock set to 5 minutes or Never during session
\item[\checkbox] 
Data collection sheet is ready
\end{itemize}

\section{13.9 Troubleshooting, YouTube Switch
Access}\label{youtube-troubleshooting}

\marginnote{Section~7 for general Switch Control and hardware troubleshooting.}

\begin{longtblr}{colspec={X[3.5,l] X[3.2,l] X[5.5,l]},rowhead=1,hlines}
\toprule
Problem & Likely Cause & Fix \\
Autoplay starts a new video automatically & Autoplay not disabled & From the video player, locate and tap the Autoplay toggle to OFF before session. Add to pre-session checklist \\
Highlight scans into YouTube home feed or search bar instead of playlist & App opened to home screen rather than playlist & Before each session, navigate to the playlist manually and leave it there. Add to pre-session checklist \\
Advertisement plays before video; unexpected scan items appear & Ad content in Switch Control scan path & Use YouTube Premium (ad-free) if available, or YouTube Kids. Use Restricted Mode to reduce ad frequency. Teacher can tap Skip Ad manually during the ad \\
Student selects wrong video accidentally & Scan path too fast or Hold Duration too short & Increase Hold Duration in Switch Control settings (Settings $>$ Accessibility $>$ Switch Control $>$ Hold Duration). Increase to 0.75-1.0 seconds \\
Video player controls appear in scan path during playback & Switch Control scanning player overlay & In early sessions, this is expected. Teach the student to scan past play controls. In later phases, this becomes a navigation skill \\
YouTube app is blocked by district MDM & District MDM policy & Contact IT to request an accessibility exception. Use YouTube Kids as an alternative. Document request in IEP \\
Student cannot return to playlist from video player & Return-to-playlist step not yet mastered & Run isolated return-to-playlist practice (Section~13.6) before resuming full sessions \\
\bottomrule
\end{longtblr}

% -- Back matter --------------------------------------------

\end{document}
